\documentclass[a4paper,12pt]{article}

% Кодировки и язык
\usepackage[T2A]{fontenc}
\usepackage[utf8]{inputenc}
\usepackage[russian,english]{babel}

% Поля
\usepackage{geometry}
\geometry{left=25mm,right=15mm,top=20mm,bottom=20mm}

% Переносы/качество верстки
\usepackage{microtype}
\emergencystretch=2em

% Таблицы и масштабирование
\usepackage{graphicx}
\usepackage{longtable}
\usepackage{float}
\usepackage{placeins}

% Выравнивание по ширине
\usepackage{ragged2e}

% Гиперссылки
\usepackage[hidelinks]{hyperref}

% Библиография (IEEE)
\usepackage[
  backend=biber,
  style=ieee,
  sorting=none
]{biblatex}
\addbibresource{references.bib}

\setlength{\parindent}{0pt}

\begin{document}
\begin{titlepage}
\thispagestyle{empty}

\begin{center}
ПРАВИТЕЛЬСТВО РОССИЙСКОЙ ФЕДЕРАЦИИ

\vspace{0.8cm}

\textbf{ФГАОУ ВО НАЦИОНАЛЬНЫЙ ИССЛЕДОВАТЕЛЬСКИЙ}\\
\textbf{УНИВЕРСИТЕТ}\\
\textbf{«ВЫСШАЯ ШКОЛА ЭКОНОМИКИ»}

\vspace{1.5cm}

Факультет компьютерных наук\\
Образовательная программа «Прикладная математика и информатика»

\vspace{4.0cm}

{\Large \textbf{Отчет о программном проекте}}

\vspace{1.8cm}

\begin{minipage}{0.92\textwidth}
\noindent на тему \textbf{Разработка чат-бота налогового консультанта}

\end{minipage}

\end{center}

\vspace{2.5cm}

\begin{flushleft}
\textbf{Выполнен студентом:}\\
Группы БПМИ246  Фадеев Павел Викторович

\vspace{2.2cm}

\textbf{Проверен руководителем проекта:}

\vspace{0.9cm}

Андреева Дарья Александровна

\vspace{0.9cm}

Старший инженер нейронных сетей

\vspace{0.9cm}

ООО «ИТ ИКС 5 ТЕХНОЛОГИИ»
\end{flushleft}

\vfill

\begin{center}
Москва 2026
\end{center}

\end{titlepage}

\renewcommand{\contentsname}{Содержание}
\tableofcontents

\newpage
\section{Введение}
\justifying
Налоговое сопровождение малого и среднего бизнеса в РФ связано с необходимостью регулярного обращения к большому объёму нормативных и разъяснительных материалов. Для корректного применения налоговых норм требуется учитывать актуальные редакции Налогового кодекса РФ, подзаконных актов, писем и разъяснений ФНС России и Минфина России, а также правоприменительную практику. При этом в реальных условиях значительная часть запросов носит прикладной характер и возникает в рамках текущих хозяйственных операций, что требует быстрых и обоснованных ответов.

Использование больших языковых моделей (LLM) в консультативных системах позволяет упростить поиск информации и представить результаты в понятной форме. Однако в задачах, где критичны точность и проверяемость, применение LLM без опоры на источники ограничено риском генерации недостоверных утверждений. В связи с этим перспективным подходом является Retrieval-Augmented Generation (RAG), при котором модель формирует ответ на основе предварительно извлечённых релевантных фрагментов документов. Такой подход повышает воспроизводимость выводов и позволяет сопровождать ответ ссылками на используемые источники.

В рамках проекта\footnote{\url{https://github.com/Pavionio/nalog_consultant_bot}} разрабатывается прототип чат-бота налогового консультанта для МСП на основе RAG-подхода. Система предназначена для обработки пользовательских запросов, извлечения релевантных фрагментов из корпуса документов и формирования ответа, содержащего обоснование и указание источников.

Практическим результатом работы является прототип системы и результаты экспериментов, позволяющие обосновать выбор ключевых компонентов архитектуры. Структура отчёта включает разделы, посвящённые постановке цели и задач, описанию архитектуры и реализации, экспериментальной оценке, анализу результатов и выводам.

\newpage
\section{Цель и задачи}

\subsection{Цель}
Разработка прототипа интеллектуального налогового консультанта для малого и среднего бизнеса, использующего методы RAG для предоставления актуальных и обоснованных консультаций.

\subsection{Задачи}
\begin{enumerate}
  \item Проанализировать предметную область налогового консультирования малого и среднего бизнеса в РФ (типовые запросы, источники норм, частые ошибки интерпретации).
  \item Исследовать существующие подходы к использованию больших языковых моделей и методов Retrieval-Augmented Generation в задачах вопросно-ответных систем и интеллектуальных консультантов.
  \item Сформировать функциональные и нефункциональные требования к системе.
  \item Разработать архитектуру RAG-системы, включая сбор, хранение, обновление и поиск данных.
  \item Реализовать прототип системы.
  \item Разработать механизм учета контекста пользовательского запроса.
  \item Провести серию экспериментов по подбору и сравнению вариантов архитектуры RAG:
    \begin{itemize}
      \item сравнить стратегии разбиения документов (chunking) и параметры перекрытия;
      \item оценить влияние выбора модели эмбеддингов и типа индекса на полноту и точность извлечения;
      \item проверить использование двухэтапного retrieval (retriever + reranker);
      \item сравнить схемы промптирования и форматы ответа (строго со ссылками, с цитированием фрагментов, с предупреждениями о неопределенности);
    \end{itemize}
  \item Провести тестирование и экспериментальную оценку качества.
  \item Проанализировать полученные результаты и оценить перспективы практического применения разработанного решения.
\end{enumerate}

\newpage
\section{Актуальность}

Налоговое регулирование характеризуется высокой динамикой и огромным количеством документов для ознакомления: регулярно вносятся изменения в Налоговый кодекс РФ, обновляются подзаконные акты, публикуются письма и разъяснения ФНС России и Минфина России, формируется и уточняется судебная практика.

Практическая потребность в оперативных и обоснованных консультациях наиболее выражена в задачах выбора налогового режима, расчёта налогов и страховых взносов, применения вычетов и льгот, подтверждения расходов, соблюдения требований к первичным документам, а также подготовки ответов на запросы и требования налоговых органов. Существенная доля затруднений обусловлена тем, что для корректного вывода требуется сопоставление нескольких источников: норм НК РФ, официальных разъяснений контролирующих органов и правоприменительной практики. Кроме того, при анализе необходимо учитывать актуальность редакций документов и применимость норм к конкретному периоду и обстоятельствам.

Использование LLM в консультативных системах потенциально позволяет повысить доступность сложной нормативной информации и сократить время на первичную обработку запросов. Чтобы поддерживать актуальность данных и контролировать галлюцинации модели целесообразно использовать Retrieval-Augmented Generation (RAG), чтобы модель формировала ответ на основе контекста, релевантного запросу.

Таким образом, разработка такого чат-бота является актуальной задачей, направленной на повышение оперативности получения справочной информации, улучшение качества интерпретации норм и снижение рисков, связанных с ошибочным применением налогового законодательства.

\newpage
\section{Обзор литературы}

\subsection{Краткая история подходов до RAG}

До появления Retrieval-Augmented Generation (RAG) в задачах вопросно-ответных систем по большим корпусам текстов доминировали два класса решений: системы «поиск $\rightarrow$ извлечение ответа», где модель читателя извлекает ответ из найденного текстового фрагмента, и генеративные модели, которые формируют ответ преимущественно из параметрической памяти, без явной опоры на внешние источники. В первом случае качество в значительной степени определяется поиском: традиционно для ранжирования документов широко применялись разреженные (sparse) методы, в частности BM25 \cite{robertson1994okapi}. Такие методы интерпретируемы и устойчивы, но зависят от лексического совпадения, что ограничивает их применимость для запросов, сформулированных «своими словами».

С развитием нейросетевых представлений текста усилился интерес к dense retrieval, где запрос и текст сопоставляются в векторном пространстве. В качестве одного из наиболее влиятельных и практичных подходов для open-domain QA закрепился Dense Passage Retrieval (DPR), использующий dual-encoder и обучаемые эмбеддинги для поиска релевантных фрагментов \cite{karpukhin2020dpr}. В этой линии retrieval становится не вспомогательным модулем, а центральным компонентом качества всей системы.

Параллельно развивались идеи интегрировать retrieval в обучение языковой модели. В работе REALM retrieval используется в процессе pre-training: модель учится обращаться к внешнему корпусу как к источнику знаний, что снижает зависимость от «хранения фактов» только в параметрах сети \cite{guu2020realm}. Важно отметить, что REALM и RAG появились почти одновременно (в рамках одного года), поэтому корректнее рассматривать REALM как важную параллельную работу, подтверждающую актуальность идеи retrieval-усиления языковых моделей, а не как однозначного «предшественника» RAG \cite{guu2020realm,lewis2020rag}.

\subsection{Первая формализация RAG: Lewis et al., 2020}

Работа Lewis et al. предложила Retrieval-Augmented Generation как общий рецепт для knowledge-intensive задач: генеративная seq2seq-модель (параметрическая память) дополняется непараметрической памятью в виде dense-индекса, из которого извлекаются релевантные текстовые фрагменты \cite{lewis2020rag}. Важный вклад этой работы состоит в том, что retrieval рассматривается как механизм обновляемости знаний без переобучения генератора и потенциальной проверяемости ответа за счёт привязки к извлечённым источникам \cite{lewis2020rag}. Авторы также рассматривают две постановки: условление генерации на одном наборе retrieved-фрагментов на всю последовательность и вариант, допускающий динамический выбор источников по токенам \cite{lewis2020rag}.

Таким образом, RAG закрепил архитектурную рамку, в которой качество системы определяется совместно качеством retrieval и способом использования извлечённого контекста в генерации, что особенно релевантно прикладным доменам, где требуются актуальность и трассируемость выводов.

\subsection{Архитектура RAG: компоненты и варианты реализации}

Практически RAG удобно описывать как два модуля: (A) индексация корпуса (offline/периодически) и (B) ответ на запрос (online). В дальнейшем RAG стал рассматриваться как модульная архитектура, в которой можно независимо улучшать retrieval, подготовку контекста и генерацию \cite{gao2023rag_survey}.

\subsubsection{Модуль индексации: корпус, chunking, эмбеддинги, индекс}

\textbf{Сбор и нормализация корпуса.} Для консультационных систем по нормативным документам критично хранить не только текст, но и метаданные: тип документа, источник, дата/редакция, идентификаторы. Это позволяет вводить фильтры при поиске и поддерживать актуальность базы без изменения генератора \cite{lewis2020rag}.

\textbf{Разбиение на фрагменты (chunking).} Поскольку retrieval обычно выполняется на уровне фрагмента, документы разбиваются на чанки фиксированной или адаптивной длины, иногда с перекрытием. Размер фрагмента и перекрытие влияют на компромисс между точностью (релевантность к вопросу) и полнотой (сохранение условий, исключений и контекста). Помимо разбиения по фиксированной длине, на практике применяется и \textit{структурное (семантическое) разбиение} --- по абзацам, пунктам, заголовкам и разделам документа, что позволяет чаще сохранять логические границы норм и условий \cite{gao2023rag_survey}. Дополнительно ограничением выступает эффективность использования длинного контекста языковыми моделями: показано, что качество может снижаться, когда релевантная информация находится «в середине» длинного контекста \cite{liu2023lost}. Поэтому стратегию chunking и сборки контекста следует рассматривать как архитектурно значимый элемент RAG.

\textbf{Эмбеддинги и dense retrieval.} Наиболее распространённый практический выбор --- bi-encoder/dual-encoder подход, где векторы запроса и фрагмента вычисляются независимо, что обеспечивает быстрый поиск по большой коллекции; DPR является каноническим примером для open-domain QA \cite{karpukhin2020dpr}. Такой retriever обычно можно заменить без изменения генератора, что соответствует модульной природе RAG \cite{gao2023rag_survey}.

\textbf{Индекс и инфраструктура поиска.} Для масштабируемого retrieval применяются методы приближённого поиска ближайших соседей и соответствующие индексы. В качестве одного из наиболее распространённых инструментов для ANN-поиска используется FAISS \cite{johnson2019faiss}. В инженерных реализациях retrieval-контур часто реализуется через специализированные векторные хранилища, которые поддерживают хранение эмбеддингов, быстрый поиск и фильтрацию по метаданным (например, Qdrant) \cite{qdrant_docs}. Выбор индекса и поддержка фильтров напрямую влияют на качество и стабильность retrieval.

\subsubsection{Онлайн-модуль: retrieval, reranking, сбор контекста, генерация}

\textbf{Шаг 1: обработка запроса и первичный retrieval.} В базовом варианте строится эмбеддинг запроса и выполняется поиск top-$k$ фрагментов. В прикладных доменах качество первичного retrieval может быть недостаточно для точного ранжирования, особенно когда несколько фрагментов тематически близки.

\textbf{Шаг 2: reranking (необязательно).} Для повышения качества распространена двухэтапная схема: быстрый retriever формирует кандидатов, затем reranker уточняет порядок. Типичный reranker реализуется как cross-encoder, совместно кодирующий пару (вопрос, фрагмент), что обычно повышает точность, но требует больше вычислений \cite{nogueira2019passage}. Для нормативных текстов это особенно важно, поскольку различия в условиях применения могут быть тонкими.

\textbf{Шаг 3: сбор контекста.} Из top-$k$ формируется контекст для промпта в LLM с ограничением на длину. Здесь важны отбор по релевантности, порядок, дедупликация и форматирование (например, нумерация источников). Ограничения длинного контекста дополнительно обосновывают необходимость аккуратной сборки контекста \cite{liu2023lost}.

\textbf{Шаг 4: генерация и политика надёжности.} Центральная идея RAG состоит в генерации ответа, обусловленной извлечёнными фрагментами \cite{lewis2020rag}. В прикладных консультативных системах критично ограничивать «вольную» генерацию: требовать атрибуцию ключевых утверждений источникам и корректный отказ при недостатке данных, поскольку галлюцинации являются системной проблемой \textit{neural text generation} \cite{ji2023hallucination}. 

Отдельно следует подчеркнуть, что идеи генерации ответа на основе \textit{нескольких} извлечённых источников развивались параллельно с RAG и получили дальнейшее развитие в подходах класса Fusion-in-Decoder (FiD), где агрегирование множества фрагментов осуществляется на стороне генератора \cite{izacard2021fid}. При этом качество ответа остаётся зависимым от retrieval и подготовки контекста, что согласуется с модульной трактовкой RAG \cite{lewis2020rag,gao2023rag_survey}.

\subsection{Развитие RAG после 2020 года}

После работы Lewis et al. исследования развивались в направлениях усиления использования нескольких источников при генерации (например, FiD) \cite{izacard2021fid}, а также в сторону более тесной интеграции retrieval и языковой модели, что было продемонстрировано в параллельных работах и последующих исследованиях retrieval-augmented обучения \cite{guu2020realm}. Одновременно сформировалась практика рассматривать RAG как модульную архитектуру, где качество определяется решениями на уровне индексации (chunking, метаданные), retrieval-контура (retriever/reranker) и политики надёжности генерации \cite{gao2023rag_survey,lewis2020rag}.

Таким образом, RAG следует рассматривать не как единый «трюк», а как архитектурный шаблон, позволяющий объединить обновляемую базу источников и генеративный ответ с проверяемостью по извлечённым фрагментам \cite{lewis2020rag}.

\newpage
\section{Описание экспериментов}

\subsection{Цель экспериментов}

Первые эксперименты были направлены на проверку работоспособности базового контура Retrieval-Augmented Generation (RAG) «сбор данных $\rightarrow$ индексация $\rightarrow$ извлечение релевантных фрагментов $\rightarrow$ генерация ответа со ссылками на источники» \cite{lewis2020rag}. Целью было получить воспроизводимый baseline, позволяющий:
\begin{itemize}
  \item выполнять семантический поиск по корпусу налоговых документов;
  \item формировать ответы, строго привязанные к извлечённому контексту;
  \item выявить первичные ограничения качества (ошибки извлечения релевантных фрагментов, недостаточность контекста, конфликт источников).
\end{itemize}

На этапе первых экспериментов использовалась одношаговая схема dense retrieval без reranking, без гибридного поиска и без автоматического обновления корпуса.

\subsection{Постановка и условия экспериментов}

В этом подразделе фиксируются общие условия, в которых сравнивались методы. Это важно, поскольку в RAG-системе результат зависит не только от выбранного retrieval-метода, но и от состава корпуса, способа разбиения документов, модели эмбеддингов, параметров генерации и критерия попадания в эталон.

\subsubsection{Экспериментальное окружение и конфигурация}

Все компоненты разворачивались локально с использованием \texttt{docker-compose}. Векторное хранилище реализовано на Qdrant, реестр документов и состояние обновления корпуса хранятся в PostgreSQL. Эксперименты запускались на локальной инфраструктуре с отдельным управлением моделями для retrieval-трансформаций, генерации ответа и LLM-валидации.

Для dense retrieval использовалась коллекция Qdrant с чанками корпуса официальных документов ФНС, Минфина и \texttt{pravo.gov.ru}. Базовая модель эмбеддингов в основной retrieval-матрице --- \texttt{BAAI/bge-m3}. Для экспериментов с query rewrite и HyDE применялась модель \texttt{openai/gpt-oss-20b} с \texttt{reasoning\_effort=low}. Такой режим выбран как компромисс: он даёт достаточно стабильные переформулировки запросов, но не превращает предварительную обработку запроса в основной вычислительный bottleneck.

Генерация ответов для последующей оценки качества генерации выполнялась отдельно от retrieval-оценки. Для этого использовалась модель \texttt{qwen/qwen3-14b} в режиме \texttt{no think}: в запрос явно добавлялся запрет на reasoning, а в OpenAI-compatible payload передавался параметр \texttt{enable\_thinking=false}, если сервер его поддерживает. Параметры генерации: \texttt{temperature=0}, \texttt{max\_tokens=1500}. Нулевая температура выбрана по трём причинам.

Во-первых, \textbf{воспроизводимость}: при оценке качества retrieval и генерации важно, чтобы при одинаковых входных данных система давала один и тот же результат. Это позволяет сравнивать конфигурации retrieval и судей без дополнительного шума от случайной генерации.

Во-вторых, \textbf{доверие пользователей в legal-tech контексте}: если система на один и тот же вопрос даёт разные ответы при каждом запуске, это подрывает доверие к ней в регуляторном домене. Пользователи ожидают, что одному набору источников соответствует один детерминированный вывод.

В-третьих, \textbf{снижение галлюцинаций}: при \texttt{temperature=0} модель выбирает наиболее вероятный токен на каждом шаге и не вносит случайные отклонения от контекста, что особенно важно при строгой политике промпта (запрет добавлять факты вне источников) \cite{ji2023hallucination}.

LLM-валидация качества генерации также вынесена в отдельный этап. Для неё предусмотрены три независимых валидатора: \texttt{precision}, \texttt{completeness} и \texttt{format}. Валидация может запускаться другой моделью, чем генерация ответа; в текущей конфигурации для LLM-судьи используется \texttt{openai/gpt-oss-20b} с \texttt{reasoning\_effort=low}. Оценка генерации не пересчитывает ответы заново, а работает по заранее сохранённому JSONL с вопросом, найденным контекстом и ответом модели.

В экспериментах с реранкером используется двухэтапная схема retrieval: быстрый dense retriever сначала возвращает расширенный набор кандидатов, после чего cross-encoder/reranker переупорядочивает их и отдаёт финальный top-$k$. Это позволяет отдельно измерять влияние первичного поиска, query transformation и reranking.

\subsubsection{Роль PostgreSQL и Qdrant в прототипе}

В прототипе используются два хранилища с разными задачами:
\begin{itemize}
  \item \textbf{PostgreSQL} применяется для ведения реестра документов и их жизненного цикла: идентификация документа по источнику и URL, хранение базовых метаданных (тип, заголовок, даты, номер), контроль обновлений и изменений (ETag/Last-Modified/хеш), планирование повторных проверок и поддержка очереди задач на загрузку/переиндексацию. Это необходимо для реализации автообновления корпуса и воспроизводимости состава данных на разных этапах экспериментов.
  \item \textbf{Qdrant} используется как векторное хранилище для retrieval: хранение эмбеддингов фрагментов, быстрый поиск top-$k$ релевантных фрагментов по эмбеддингу запроса, а также хранение payload-метаданных, необходимых для цитирования (например, URL и идентификаторы). Такой раздел обязанностей позволяет независимо развивать контур обновления данных и контур семантического поиска.
\end{itemize}

\subsubsection{Источники данных и текущий охват корпуса}

В конфигурации источников определены 15 активных ресурсов трёх категорий:

\textbf{ФНС России:}
\begin{itemize}
  \item \texttt{nalog\_letters} --- письма и разъяснения ФНС (\url{https://www.nalog.gov.ru/rn77/about_fts/about_nalog/});
  \item \texttt{nalog\_docs} --- нормативные документы ФНС (\url{https://www.nalog.gov.ru/rn77/about_fts/docs/});
  \item \texttt{nalog\_calendar} --- налоговый календарь ФНС (XML).
\end{itemize}

\textbf{Минфин России} (раздел «Ответы на вопросы», разбитый по тематическим подразделам):
\texttt{minfin\_commonlaw}, \texttt{minfin\_orgprofit}, \texttt{minfin\_fizprofit}, \texttt{minfin\_property}, \texttt{minfin\_indirect}, \texttt{minfin\_international}, \texttt{minfin\_special}, \texttt{minfin\_transfert}, \texttt{minfin\_foreign}, \texttt{minfin\_customs\_value}, \texttt{minfin\_imposition}.

\textbf{Законодательство:}
\begin{itemize}
  \item \texttt{pravo\_nk1} --- Налоговый кодекс РФ (\url{http://pravo.gov.ru}).
\end{itemize}

Данные извлекаются автоматически: для каждого источника реализован обработчик (\texttt{handler}), выполняющий обход страниц, извлечение текста, чанкинг и индексацию в Qdrant. Периодичность обновления задаётся параметром \texttt{crawl\_freq\_days} для каждого источника (как правило, 7 дней).

\subsubsection{Хранение метаданных и автообновление корпуса}

Для контроля изменений и поддержания актуальности базы знаний реализована схема хранения метаданных в PostgreSQL. Таблица \texttt{rag\_doc} хранит карточку каждого документа: идентификатор источника, URL, метаданные, контрольную сумму содержимого и HTTP-валидаторы (ETag, Last-Modified), а также время следующей плановой проверки (\texttt{next\_check\_at}). Таблица \texttt{rag\_fetch\_log} ведёт журнал всех обращений к источникам.

Автообновление запускается командой и работает следующим образом: из \texttt{rag\_doc} выбираются документы, у которых \texttt{next\_check\_at} наступило, выполняется повторная загрузка, и если хеш или HTTP-валидаторы изменились, документ переиндексируется в Qdrant. Это позволяет отслеживать обновления нормативной базы без повторного полного обхода источников.

\subsubsection{Индексация и retrieval}

Для dense retrieval применялась модель эмбеддингов \texttt{BAAI/bge-m3} \cite{m3embedding2024}. Эмбеддинги вычислялись через \texttt{SentenceTransformer} с нормализацией (\texttt{normalize\_embeddings=True}). Поиск реализован как запрос к Qdrant с возвратом \texttt{top\_k} наиболее релевантных фрагментов и payload-метаданных.

Для retrieval использовались параметры:
\begin{itemize}
  \item \texttt{top\_k=1}, \texttt{top\_k=5} и \texttt{top\_k=10} (сравнение в экспериментах);
  \item \texttt{score\_threshold=0.0} (без отсечения по порогу на этапе baseline).
\end{itemize}

\subsubsection{Разбиение документов на фрагменты и сбор контекста}

В baseline-реализации документы представлены в виде набора текстовых фрагментов. На этапе сборки контекста применялись ограничения \texttt{max\_chunk\_chars=2500} (обрезка одного фрагмента) и \texttt{max\_context\_chars=18000} (общий лимит контекста). Фрагменты добавлялись в контекст в порядке выдачи retrieval до достижения лимита; каждый фрагмент нумеровался как \texttt{[1]}, \texttt{[2]}, \dots, что далее используется для атрибуции в ответе.

\subsubsection{Генерация ответа и политика надёжности}

Генерация ответа отделена от retrieval-оценки. Retrieval-эксперименты измеряют, попал ли нужный документ или чанк в выдачу; генерационная оценка проверяет, насколько модель корректно использовала уже найденный контекст. Для генерации ответов использовалась модель \texttt{qwen/qwen3-14b} в режиме \texttt{no think}.

Основной вариант промпта генерации --- \texttt{strict\_citations}. Он задаёт модель как налогового консультанта, запрещает использовать внешние знания и требует отвечать только по предоставленному контексту. Если в контексте нет достаточной информации, модель должна явно отказаться: «В предоставленных документах нет достаточной информации для ответа». Ответ должен иметь структуру: краткий вывод, обоснование по документам и источники. Также промпт отдельно запрещает выдумывать реквизиты, ссылки, сроки, суммы штрафов и номера статей.

Такая политика мотивирована проблемой галлюцинаций в neural text generation \cite{ji2023hallucination} и соответствует общей идее RAG как способа повысить проверяемость ответа \cite{lewis2020rag}. Важно, что отказ при недостатке контекста считается корректным поведением: модель не должна достраивать налоговый или правовой вывод из собственных знаний, если найденные документы не дают основания для ответа.

\subsubsection{«Золотой набор» для оценки качества retrieval}

\paragraph{Принцип формирования.}

Для количественной оценки retrieval был сформирован размеченный набор запросов~--- «золотой набор» (eval dataset). Каждый элемент набора содержит вопрос и список эталонных релевантных объектов. Поле \texttt{query} содержит пользовательский вопрос, поле \texttt{relevant} --- один или несколько релевантных документов или чанков с \texttt{source\_code}, \texttt{external\_id} и, в новом формате, \texttt{chunk\_i}. Важно, что для одного вопроса может быть указано несколько релевантных чанков: это отражает ситуацию, когда ответ подтверждается несколькими фрагментами одного документа или несколькими документами. Такая разметка позволяет считать как попадание на уровне документа, так и более строгую оценку на уровне чанка.

Набор формировался полностью автоматически скриптом \texttt{gen\_eval\_dataset.py} в несколько шагов:
\begin{enumerate}
  \item Из Qdrant извлекаются все чанки с текстом.
  \item Выполняется стратифицированная выборка: до $n=18$ документов на каждый \texttt{source\_code}, дедупликация по \texttt{external\_id} (для одного документа берётся первый чанк), фиксированный \texttt{seed=42}.
  \item Для каждого чанка делается запрос к LLM: промпт задаёт уровень сложности и требует вернуть ровно один вопрос.
  \item Ответ фильтруется: отбрасываются вопросы короче 20 символов и содержащие артефакты промпта («фрагмент», «приведён», «из текста»).
\end{enumerate}
Итоговый набор сохраняется в формате JSONL. Дополнительно в новом формате разметки поддерживается поле \texttt{hard\_negatives}: оно содержит похожие, но нерелевантные документы или чанки. Такие примеры нужны для анализа ошибок, когда система находит тематически близкий фрагмент, но он не отвечает на исходный вопрос из-за отличий в налоге, периоде, категории налогоплательщика или условии применения нормы.

Для оценки качества генерации строится отдельный JSONL-набор. Он создаётся по уже существующим eval dataset: для каждого вопроса выполняется retrieval выбранным методом, найденные чанки сопоставляются с \texttt{relevant} и \texttt{hard\_negatives}, затем по найденному контексту генерируется ответ. В сохранённой записи фиксируются вопрос, ответ, найденные чанки, метки \texttt{gold\_relevant}/\texttt{hard\_negative}/\texttt{unknown}, краткую сводку качества найденного контекста и поля ручной оценки. По умолчанию ручные оценки равны \texttt{-1}; это означает, что человек ещё не оценил пример. Такие значения не участвуют в расчёте согласованности человека и LLM-судьи, но не мешают посчитать и сохранить оценку LLM-валидаторов.

\paragraph{Три уровня сложности.}

Чтобы покрыть разные сценарии использования системы~--- от профессионального юриста до рядового предпринимателя~--- было сформировано три датасета разного уровня сложности. Уровень сложности задаётся промптом при генерации вопроса.

\textbf{Easy.} Вопросы генерировались с промптом, требующим «конкретный вопрос реального пользователя» без дополнительных ограничений по стилю. На выходе получились формальные, юридически точные вопросы: все начинаются с заглавной буквы, используют официальную лексику и часто ссылаются на конкретные нормы, законы, статьи и документы (медиана длины~--- 163 символа, аббревиатуры в 38\% вопросов). Такой стиль ближе к запросу профессионального юриста, чем рядового предпринимателя. Пример:

\begin{quote}
\textit{«Какие федеральные законы изменили правила налогообложения доходов физических лиц, превышающих 5 млн рублей за налоговый период, с 1 января 2021 года?»}
\end{quote}

\textbf{Hard.} Промпт требовал сформулировать вопрос косвенно, разговорным языком, с опечатками. На выходе вопросы заметно короче (медиана~--- 84 символа) и ближе к реальному пользователю: 57\% начинаются со строчной буквы, встречаются разговорные обороты и опечатки. Аббревиатуры при этом сохраняются (51\% вопросов: НДС, УСН, НКО и т.п.). Пример:

\begin{quote}
\textit{«мне надо плачить актиз за этиловый спирт? какая ставка там будет?»}
\end{quote}

\textbf{Superhard.} Самый сложный уровень: вопрос с обильным бизнес-контекстом, опечатками, разговорными оборотами и смещённой формулировкой. Промпт требовал каждый раз придумывать новый тип бизнеса, регион и ситуацию. Это самые длинные и зашумлённые вопросы (медиана~--- 265 символов, до 678; 82\% со строчной буквы, 16\% без знака вопроса, аббревиатуры в 75\%). Пример:

\begin{quote}
\textit{«я нашёл письмо ФНС о акцизах, но у нас в КРПИ (корпусный ремонтный парк) делаем штукатурку и штука 3-4 штуки спирта для смазки инструментов, надо ли нам платить акциз по нему, когда?»}
\end{quote}

\paragraph{Характеристики датасетов.}

Каждый из трёх датасетов содержит 171 вопрос. Их сравнительные характеристики приведены в таблице~\ref{tab:dataset_stats}.

\begin{table}[H]
\centering
\caption{Сравнение датасетов по уровням сложности}
\label{tab:dataset_stats}
\begin{tabular}{lccc}
\hline
\textbf{Параметр} & \textbf{Easy} & \textbf{Hard} & \textbf{Superhard} \\
\hline
Всего вопросов & 171 & 171 & 171 \\
Уникальных & 171 (100\%) & 171 (100\%) & 171 (100\%) \\
\hline
\multicolumn{4}{l}{\textit{Длина вопросов (символов)}} \\
минимум / медиана / максимум & 79 / 163 / 426 & 42 / 84 / 152 & 139 / 265 / 678 \\
90-й перцентиль & 223 & 117 & 401 \\
\hline
\multicolumn{4}{l}{\textit{Лексика и стиль}} \\
Со строчной буквы (разговорный стиль) & 0 (0\%) & 97 (57\%) & 141 (82\%) \\
Содержат аббревиатуры & 65 (38\%) & 87 (51\%) & 129 (75\%) \\
Без знака вопроса в конце & 0 (0\%) & 2 (1\%) & 27 (16\%) \\
TTR (type-token ratio) & 0,366 & 0,357 & 0,247 \\
\hline
\end{tabular}
\end{table}
\FloatBarrier

\textbf{Лингвистический анализ по уровням.}

\textit{Easy} демонстрирует формальный официальный стиль: 100\% вопросов начинаются с заглавной буквы, используют юридическую терминологию и часто ссылаются на конкретные нормы, законы и документы. Вопросы длинные (медиана 163 символа) и точные; аббревиатуры встречаются в 38\% вопросов. Такой стиль соответствует скорее запросу профессионального юриста, нежели рядового предпринимателя.

\textit{Hard} приближён к реальному пользователю по стилю: больше половины вопросов (57\%) написаны со строчной буквы, они заметно короче (медиана 84 символа), встречаются разговорные обороты и опечатки (например, «плачить актиз»). При этом аббревиатуры сохраняются (51\%: НДС, УСН, НКО и т.п.). Лексически hard разнообразен (TTR~0,357) благодаря коротким вопросам с разными словами.

\textit{Superhard} реалистично имитирует неструктурированные запросы: 82\% вопросов написаны со строчной буквы, 16\% не заканчиваются знаком вопроса, текст насыщен бизнес-контекстом, опечатками и разговорными оборотами. Это самые длинные вопросы (медиана 265, до 678 символов) с наибольшей долей аббревиатур (75\%). Несмотря на большой лексический объём (2161 уникальное слово), TTR невысок (0,247) из-за повторяющихся конструкций («у нас…», «надо ли платить…»).

\paragraph{Критерий попадания.}

В основной retrieval-оценке используется строгий критерий попадания на уровне чанка (\texttt{match\_mode=strict}): результат считается релевантным только если среди top-$k$ найден фрагмент, у которого одновременно совпадают \texttt{source\_code}, \texttt{external\_id} и \texttt{chunk\_i} с одним из объектов в поле \texttt{relevant}. Именно по этому критерию вычислены все основные метрики (Precision@$k$, Recall@$k$/Hit@$k$, MRR, nDCG) в сводных таблицах. Дополнительно, как вспомогательные диагностические колонки, считаются более мягкие варианты: \texttt{doc} --- попадание на уровне документа (совпадают только \texttt{source\_code} и \texttt{external\_id}, без учёта номера чанка) и \texttt{doc\_overlap} --- попадание по текстовому пересечению фрагментов. Они помогают отличить случай «найден правильный документ, но не тот чанк» от полного промаха, но в основные таблицы не выносятся.

Если у вопроса указано несколько релевантных объектов, попаданием считается нахождение любого из них. Это соответствует практическому сценарию RAG: для корректного ответа системе достаточно найти один из фрагментов, содержащих необходимое основание, хотя для полноты генерации может быть полезно найти несколько релевантных чанков.

Рассчитываются следующие метрики при $k \in \{1, 5, 10\}$:
\begin{itemize}
  \item \textbf{Hit Rate@k} --- доля запросов, для которых хотя бы один релевантный документ или релевантный чанк присутствует среди top-$k$ результатов;
  \item \textbf{MRR@k} (Mean Reciprocal Rank) --- среднее обратное ранга первого релевантного результата \cite{voorhees1999mrr};
  \item \textbf{NDCG@k} (Normalized Discounted Cumulative Gain) --- метрика качества ранжирования, учитывающая позицию релевантных результатов в выдаче;
  \item \textbf{Precision@k} --- доля релевантных результатов среди top-$k$;
  \item \textbf{HardNeg@k} --- доля запросов, для которых среди top-$k$ найден хотя бы один hard negative: тематически похожий, но нерелевантный документ или чанк. Для налогового домена эта метрика важна, поскольку похожие документы могут отличаться налогом, периодом, категорией налогоплательщика или условием применения нормы.
\end{itemize}

\subsubsection{Сводка экспериментальных настроек}

В таблице~\ref{tab:experiment_settings} приведены фиксированные параметры для основных групп экспериментов. В каждой группе менялся только тот компонент, который указан в столбце «Варьируемый фактор»; остальные параметры оставались фиксированными. Это позволяет интерпретировать различия в метриках как эффект конкретного изменения, а не как следствие одновременной смены нескольких частей пайплайна.

\begin{table}[H]
\centering
\caption{Сводка условий проведения экспериментов}
\label{tab:experiment_settings}
\resizebox{\textwidth}{!}{%
\begin{tabular}{lll}
\hline
\textbf{Группа экспериментов} & \textbf{Варьируемый фактор} & \textbf{Фиксированные параметры} \\
\hline
Retrieval matrix & Метод retrieval & Qdrant, \texttt{BAAI/bge-m3}, $k \in \{1,5,10\}$, \texttt{score\_threshold=0} \\
Rewrite / HyDE & Трансформация запроса & \texttt{openai/gpt-oss-20b}, \texttt{reasoning\_effort=low}, тот же dense retriever \\
Preprocessing & Профиль очистки текста & baseline dense retrieval, \texttt{BAAI/bge-m3}, $k=5$ \\
Embedder & Модель эмбеддингов & baseline retrieval, одинаковый корпус, $k=5$, без reranker \\
Chunking & Стратегия разбиения & baseline dense retrieval, $k=5$, сопоставление document/overlap \\
Reranker & Модель и \texttt{fetch\_k} & dense candidate retrieval, затем reranking до \texttt{final\_k} \\
Генерация ответа & Модель генерации и prompt & \texttt{qwen/qwen3-14b}, \texttt{no think}, \texttt{temperature=0}, \texttt{strict\_citations} \\
LLM-судья & Модель и валидаторы & \texttt{openai/gpt-oss-20b}, \texttt{reasoning\_effort=low}, \texttt{precision/completeness/format} \\
\hline
\end{tabular}%
}
\end{table}
\FloatBarrier

\subsection{Наблюдения по результатам первичных прогонов}

Первые прогоны выполнялись вручную через CLI-интерфейс. По итогам были получены следующие наблюдения:
\begin{itemize}
  \item \textbf{Работоспособность сквозного контура.} Система стабильно извлекает релевантные фрагменты из Qdrant и формирует ответы, сопровождая их ссылками на источники (по нумерации \texttt{[n]}).
  \item \textbf{Зависимость качества от retrieval.} Основной класс ошибок baseline связан с тем, что среди top-$k$ отсутствует фрагмент, содержащий нужное условие, исключение или определение. В таких случаях модель корректно отказывает или запрашивает недостающий контекст (в соответствии с промптом), но полезность ответа ограничивается качеством retrieval.
  \item \textbf{Чувствительность к разбиению.} При неудачном разбиении часть юридически значимых условий может «разрезаться» между фрагментами, что повышает долю ситуаций «недостаточно данных в контексте» даже при наличии документа в корпусе. Это подтвердило необходимость отдельной серии экспериментов с chunking и перекрытием.
  \item \textbf{Роль строгого формата ответа.} Требование ссылок \texttt{[n]} и отказа при недостатке данных снижает долю недостоверных утверждений, но повышает требования к полноте retrieval (иначе система чаще отказывает).
\end{itemize}

\subsection{Результаты экспериментов по retrieval}

Была запущена полная матрица retrieval-методов на трёх датасетах при $k \in \{1,5,10\}$. Сравнивались пять конфигураций: baseline dense retrieval, query rewrite, HyDE, reranker и HyDE+reranker. Baseline выполняет прямой dense-поиск по эмбеддингу исходного вопроса. Rewrite сначала переписывает пользовательский запрос в более формальную налоговую формулировку, HyDE генерирует гипотетический фрагмент документа и ищет уже по нему, reranker переупорядочивает найденных dense-кандидатов, а HyDE+reranker объединяет гипотетический документ и второй этап ранжирования. Для rewrite и HyDE использовалась модель \texttt{openai/gpt-oss-20b} с \texttt{reasoning\_effort=low}. Сначала приведены агрегированные результаты по всем трём датасетам, затем --- разрезы по уровням сложности.

\clearpage
\begin{table}[H]
\centering
\caption{Агрегированные retrieval-метрики по трём датасетам}
\label{tab:retrieval_overall}
\resizebox{\textwidth}{!}{%
\begin{tabular}{llccccc}
\hline
\textbf{Метод} & \textbf{k} & \textbf{Hit@k} & \textbf{MRR@k} & \textbf{nDCG@k} & \textbf{Precision@k} & \textbf{HardNeg@k} \\
\hline
baseline & 1 & 0,602 & 0,602 & 0,602 & 0,602 & 0,166 \\
rewrite & 1 & 0,526 & 0,526 & 0,526 & 0,526 & 0,125 \\
HyDE & 1 & 0,343 & 0,343 & 0,343 & 0,343 & 0,107 \\
reranker & 1 & \textbf{0,708} & \textbf{0,708} & \textbf{0,708} & \textbf{0,708} & 0,121 \\
HyDE+reranker & 1 & 0,429 & 0,429 & 0,429 & 0,429 & 0,125 \\
\hline
baseline & 5 & 0,789 & 0,673 & 0,699 & 0,179 & 0,343 \\
rewrite & 5 & 0,737 & 0,608 & 0,639 & 0,163 & 0,287 \\
HyDE & 5 & 0,528 & 0,412 & 0,439 & 0,117 & 0,230 \\
reranker & 5 & \textbf{0,836} & \textbf{0,761} & \textbf{0,778} & \textbf{0,186} & 0,327 \\
HyDE+reranker & 5 & 0,593 & 0,496 & 0,520 & 0,130 & 0,246 \\
\hline
baseline & 10 & 0,799 & 0,674 & 0,703 & 0,151 & 0,345 \\
rewrite & 10 & 0,749 & 0,610 & 0,643 & 0,138 & 0,294 \\
HyDE & 10 & 0,552 & 0,416 & 0,448 & 0,101 & 0,238 \\
reranker & 10 & \textbf{0,836} & \textbf{0,761} & \textbf{0,779} & \textbf{0,157} & 0,333 \\
HyDE+reranker & 10 & 0,604 & 0,498 & 0,525 & 0,111 & 0,250 \\
\hline
\end{tabular}%
}
\end{table}
\FloatBarrier

\begin{table}[H]
\centering
\caption{Retrieval-метрики на easy dataset}
\label{tab:retrieval_easy}
\resizebox{\textwidth}{!}{%
\begin{tabular}{llccccc}
\hline
\textbf{Метод} & \textbf{k} & \textbf{Hit@k} & \textbf{MRR@k} & \textbf{nDCG@k} & \textbf{Precision@k} & \textbf{HardNeg@k} \\
\hline
baseline & 1  & 0,795 & 0,795 & 0,795 & 0,795 & 0,088 \\
rewrite & 1  & 0,772 & 0,772 & 0,772 & 0,772 & 0,053 \\
HyDE & 1  & 0,561 & 0,561 & 0,561 & 0,561 & 0,088 \\
reranker & 1  & \textbf{0,871} & \textbf{0,871} & \textbf{0,871} & \textbf{0,871} & 0,082 \\
HyDE+reranker & 1  & 0,696 & 0,696 & 0,696 & 0,696 & 0,070 \\
\hline
baseline & 5  & 0,924 & 0,850 & 0,867 & 0,195 & 0,275 \\
rewrite & 5  & 0,924 & 0,835 & 0,858 & 0,193 & 0,257 \\
HyDE & 5  & 0,754 & 0,636 & 0,665 & 0,158 & 0,216 \\
reranker & 5  & \textbf{0,953} & \textbf{0,908} & \textbf{0,917} & \textbf{0,199} & 0,257 \\
HyDE+reranker & 5  & 0,830 & 0,750 & 0,769 & 0,173 & 0,228 \\
\hline
baseline & 10 & 0,936 & 0,852 & 0,873 & 0,166 & 0,275 \\
rewrite & 10 & 0,930 & 0,836 & 0,860 & 0,162 & 0,257 \\
HyDE & 10 & 0,766 & 0,638 & 0,670 & 0,135 & 0,216 \\
reranker & 10 & \textbf{0,953} & \textbf{0,908} & \textbf{0,918} & \textbf{0,167} & 0,269 \\
HyDE+reranker & 10 & 0,836 & 0,751 & 0,771 & 0,145 & 0,234 \\
\hline
\end{tabular}%
}
\end{table}
\FloatBarrier

\begin{table}[H]
\centering
\caption{Retrieval-метрики на hard dataset}
\label{tab:retrieval_hard}
\resizebox{\textwidth}{!}{%
\begin{tabular}{llccccc}
\hline
\textbf{Метод} & \textbf{k} & \textbf{Hit@k} & \textbf{MRR@k} & \textbf{nDCG@k} & \textbf{Precision@k} & \textbf{HardNeg@k} \\
\hline
baseline & 1  & 0,456 & 0,456 & 0,456 & 0,456 & 0,228 \\
rewrite & 1  & 0,404 & 0,404 & 0,404 & 0,404 & 0,170 \\
HyDE & 1  & 0,246 & 0,246 & 0,246 & 0,246 & 0,129 \\
reranker & 1  & \textbf{0,526} & \textbf{0,526} & \textbf{0,526} & \textbf{0,526} & 0,164 \\
HyDE+reranker & 1  & 0,281 & 0,281 & 0,281 & 0,281 & 0,181 \\
\hline
baseline & 5  & 0,637 & 0,521 & 0,548 & 0,152 & 0,345 \\
rewrite & 5  & 0,626 & 0,490 & 0,520 & 0,144 & 0,292 \\
HyDE & 5  & 0,444 & 0,320 & 0,349 & 0,102 & 0,240 \\
reranker & 5  & \textbf{0,702} & \textbf{0,598} & \textbf{0,622} & \textbf{0,163} & 0,345 \\
HyDE+reranker & 5  & 0,480 & 0,364 & 0,392 & 0,113 & 0,263 \\
\hline
baseline & 10 & 0,643 & 0,522 & 0,550 & 0,128 & 0,351 \\
rewrite & 10 & 0,643 & 0,493 & 0,527 & 0,124 & 0,304 \\
HyDE & 10 & 0,444 & 0,320 & 0,349 & 0,085 & 0,251 \\
reranker & 10 & \textbf{0,702} & \textbf{0,598} & \textbf{0,624} & \textbf{0,137} & 0,345 \\
HyDE+reranker & 10 & 0,503 & 0,368 & 0,402 & 0,100 & 0,263 \\
\hline
\end{tabular}%
}
\end{table}
\FloatBarrier

\begin{table}[H]
\centering
\caption{Retrieval-метрики на superhard dataset}
\label{tab:retrieval_superhard}
\resizebox{\textwidth}{!}{%
\begin{tabular}{llccccc}
\hline
\textbf{Метод} & \textbf{k} & \textbf{Hit@k} & \textbf{MRR@k} & \textbf{nDCG@k} & \textbf{Precision@k} & \textbf{HardNeg@k} \\
\hline
baseline & 1  & 0,556 & 0,556 & 0,556 & 0,556 & 0,181 \\
rewrite & 1  & 0,404 & 0,404 & 0,404 & 0,404 & 0,152 \\
HyDE & 1  & 0,222 & 0,222 & 0,222 & 0,222 & 0,105 \\
reranker & 1  & \textbf{0,725} & \textbf{0,725} & \textbf{0,725} & \textbf{0,725} & 0,117 \\
HyDE+reranker & 1  & 0,310 & 0,310 & 0,310 & 0,310 & 0,123 \\
\hline
baseline & 5  & 0,807 & 0,647 & 0,682 & 0,188 & 0,409 \\
rewrite & 5  & 0,661 & 0,499 & 0,539 & 0,152 & 0,310 \\
HyDE & 5  & 0,386 & 0,279 & 0,304 & 0,090 & 0,234 \\
reranker & 5  & \textbf{0,854} & \textbf{0,778} & \textbf{0,794} & \textbf{0,198} & 0,380 \\
HyDE+reranker & 5  & 0,468 & 0,374 & 0,397 & 0,104 & 0,246 \\
\hline
baseline & 10 & 0,819 & 0,649 & 0,687 & 0,160 & 0,409 \\
rewrite & 10 & 0,673 & 0,501 & 0,543 & 0,129 & 0,322 \\
HyDE & 10 & 0,444 & 0,289 & 0,325 & 0,085 & 0,246 \\
reranker & 10 & \textbf{0,854} & \textbf{0,778} & \textbf{0,795} & \textbf{0,166} & 0,386 \\
HyDE+reranker & 10 & 0,474 & 0,375 & 0,400 & 0,089 & 0,251 \\
\hline
\end{tabular}%
}
\end{table}
\FloatBarrier

Ключевые наблюдения по результатам:
\begin{itemize}
  \item \textbf{Reranker} остаётся лучшей конфигурацией на всех трёх датасетах и всех значениях $k$. На hard при $k=5$ Hit@5 вырастает с 0,637 у baseline до 0,702, а на superhard --- с 0,807 до 0,854. Наиболее заметный эффект реранкера проявляется в MRR: на superhard при $k=5$ MRR растёт с 0,647 до 0,778, то есть релевантный документ чаще поднимается на первые позиции.
  \item \textbf{Query rewrite} не даёт устойчивого улучшения. На easy при $k=5$ Hit@5 совпадает с baseline (0,924), но MRR ниже. На hard и superhard rewrite хуже baseline, что указывает на потерю части полезных пользовательских формулировок при переводе запроса в более формальный налоговый язык.
  \item \textbf{HyDE} показывает худшие результаты среди основных конфигураций. Особенно сильное падение видно на hard и superhard: при $k=5$ Hit@5 равен 0,444 и 0,386 соответственно. Это означает, что гипотетический документ часто уводит dense retrieval от исходного релевантного документа.
  \item \textbf{HyDE+reranker} улучшает чистый HyDE, но не догоняет baseline. Это подтверждает ограничение двухэтапной схемы: reranker может переупорядочить найденных кандидатов, но не может восстановить релевантный документ, если он не попал в набор кандидатов на этапе первичного поиска.
  \item Увеличение $k$ с 5 до 10 даёт небольшой прирост Hit@k, но почти не меняет MRR. Это означает, что дополнительные релевантные документы чаще появляются на нижних позициях, а качество первых результатов определяется прежде всего reranking.
  \item \textbf{Hard negatives остаются значимым источником ошибок.} Даже при росте Hit@k часть выдачи состоит из тематически близких, но нерелевантных документов. Это особенно заметно в налоговом домене: фрагменты могут содержать те же налоги или термины, но относиться к другому периоду, другому типу налогоплательщика или другому условию применения нормы.
\end{itemize}

\subsection{Эксперименты с предобработкой текста}

Отдельно проверялось влияние очистки документов перед индексацией. Мотивация этого эксперимента состоит в том, что HTML-страницы Минфина и ФНС часто содержат навигационные блоки, повторяющиеся заголовки, меню, служебный текст и элементы интерфейса. Такие фрагменты попадают в чанки, размывают семантическое представление документа и могут приводить к ложным совпадениям по словам, не относящимся к правовой норме.

Важное замечание о масштабе метрик. Эта серия и две следующие (эмбеддеры, разбиение на чанки) выполнялись на отдельно переиндексированном локальном корпусе и с более мягким критерием попадания \texttt{doc\_overlap}. Поэтому абсолютные значения Hit@5 здесь заметно ниже, чем в основной retrieval-матрице (таблицы~\ref{tab:retrieval_easy}--\ref{tab:retrieval_superhard}), и сравнивать их корректно только \emph{внутри} серии, а не с основными таблицами (подробнее см. раздел~\ref{sec:optimal}).

Во всех вариантах использовался baseline dense retrieval при $k=5$; менялся только профиль предобработки текста. Профиль \texttt{raw} сохраняет исходный очищенный текст без дополнительного удаления boilerplate. \texttt{clean\_basic} и \texttt{clean\_legal} применяют более осторожные правила нормализации. \texttt{clean\_aggressive} удаляет больше повторяющихся элементов, но повышает риск потерять полезный контекст. \texttt{clean\_no\_boilerplate} наиболее явно нацелен на удаление навигационного шума.

Результаты показывают, что удаление boilerplate может существенно повысить Hit@5 относительно raw-представления: профиль \texttt{clean\_no\_boilerplate} показывает лучшие значения Hit@5, MRR@5 и nDCG@5 на всех трёх датасетах. Однако этот профиль одновременно увеличивает HardNeg@5. Это означает, что после удаления навигационного шума документы становятся более «плотными» по юридической терминологии, и retriever чаще поднимает тематически похожие, но нерелевантные фрагменты. Поэтому \texttt{clean\_no\_boilerplate} выглядит перспективным кандидатом, но требует проверки совместно с реранкером и ручного просмотра очищенных документов.

\begin{table}[H]
\centering
\caption{Сравнение профилей предобработки при $k=5$}
\label{tab:preprocessing_results}
\resizebox{\textwidth}{!}{%
\begin{tabular}{llccccc}
\hline
\textbf{Профиль} & \textbf{Датасет} & \textbf{Hit@5} & \textbf{MRR@5} & \textbf{nDCG@5} & \textbf{Precision@5} & \textbf{HardNeg@5} \\
\hline
raw & Easy & 0,292 & 0,181 & 0,209 & 0,063 & 0,149 \\
raw & Hard & 0,140 & 0,096 & 0,107 & 0,035 & 0,210 \\
raw & Superhard & 0,193 & 0,140 & 0,153 & 0,044 & 0,230 \\
clean\_basic & Easy & 0,287 & 0,181 & 0,207 & 0,062 & 0,170 \\
clean\_basic & Hard & 0,140 & 0,094 & 0,106 & 0,035 & 0,194 \\
clean\_basic & Superhard & 0,193 & 0,141 & 0,154 & 0,044 & 0,243 \\
clean\_legal & Easy & 0,275 & 0,173 & 0,199 & 0,060 & 0,191 \\
clean\_legal & Hard & 0,140 & 0,097 & 0,106 & 0,035 & 0,210 \\
clean\_legal & Superhard & 0,193 & 0,140 & 0,154 & 0,046 & 0,230 \\
clean\_aggressive & Easy & 0,281 & 0,195 & 0,217 & 0,061 & 0,277 \\
clean\_aggressive & Hard & 0,158 & 0,117 & 0,125 & 0,041 & 0,242 \\
clean\_aggressive & Superhard & 0,228 & 0,161 & 0,178 & 0,053 & 0,351 \\
clean\_no\_boilerplate & Easy & \textbf{0,404} & \textbf{0,266} & \textbf{0,298} & \textbf{0,085} & 0,277 \\
clean\_no\_boilerplate & Hard & \textbf{0,240} & \textbf{0,145} & \textbf{0,168} & \textbf{0,055} & 0,323 \\
clean\_no\_boilerplate & Superhard & \textbf{0,287} & \textbf{0,217} & \textbf{0,235} & \textbf{0,064} & 0,338 \\
\hline
\end{tabular}%
}
\end{table}
\FloatBarrier

\textbf{Вывод.} Для боевого пайплайна выбран профиль \texttt{clean\_no\_boilerplate} как дающий наибольший прирост Hit@5, MRR@5 и nDCG@5 на всех трёх уровнях сложности. Сопутствующий рост HardNeg@5 не критичен: на следующем этапе кандидатов переупорядочивает реранкер, который отсеивает тематически близкие, но нерелевантные фрагменты.

\subsection{Эксперименты с эмбеддерами}

Сравнение эмбеддеров выполнялось при одинаковом baseline retrieval и $k=5$: корпус, критерий попадания, Qdrant-поиск и отсутствие reranker оставались неизменными, менялась только модель построения эмбеддингов. Цель эксперимента --- отделить качество семантического представления текста от остальных факторов пайплайна.

Для налоговых документов это особенно важно, потому что запросы могут быть сформулированы разговорно, а документы написаны формальным юридическим языком. Хороший эмбеддер должен сопоставлять такие формулировки без буквального совпадения терминов, но при этом не смешивать похожие налоговые ситуации с разными условиями применения. В итоговую таблицу включены только модели, прошедшие корректный smoke-тест и использовавшие валидный режим кодирования запроса и документа. Модели, для которых в ходе проверки были выявлены проблемы загрузки, некорректный формат embedding-инференса или невалидная конфигурация, исключены из итогового сравнения.

Лучший результат по Hit@5 на всех трёх датасетах показал \texttt{e5\_large}; при этом \texttt{e5\_base} оказался быстрее и заметно лучше \texttt{BAAI/bge-m3} на hard dataset.

\begin{table}[H]
\centering
\caption{Сравнение эмбеддеров при $k=5$}
\label{tab:embedder_results}
\resizebox{\textwidth}{!}{%
\begin{tabular}{llccccc}
\hline
\textbf{Эмбеддер} & \textbf{Датасет} & \textbf{Hit@5} & \textbf{MRR@5} & \textbf{nDCG@5} & \textbf{Precision@5} & \textbf{p95, мс} \\
\hline
e5\_large & Easy & \textbf{0,327} & \textbf{0,209} & \textbf{0,239} & \textbf{0,070} & 59,9 \\
bge\_m3 & Easy & 0,292 & 0,180 & 0,208 & 0,063 & 50,5 \\
e5\_base & Easy & 0,281 & 0,174 & 0,200 & 0,061 & \textbf{33,3} \\
jina\_embeddings\_v3 & Easy & 0,263 & 0,164 & 0,189 & 0,057 & 200,3 \\
\hline
e5\_large & Hard & \textbf{0,216} & \textbf{0,135} & \textbf{0,155} & \textbf{0,049} & 63,6 \\
e5\_base & Hard & 0,193 & 0,120 & 0,139 & 0,047 & \textbf{30,0} \\
jina\_embeddings\_v3 & Hard & 0,164 & 0,104 & 0,119 & 0,041 & 206,1 \\
bge\_m3 & Hard & 0,135 & 0,093 & 0,103 & 0,034 & 51,1 \\
\hline
e5\_large & Superhard & \textbf{0,246} & \textbf{0,164} & \textbf{0,184} & \textbf{0,054} & 65,4 \\
e5\_base & Superhard & 0,199 & 0,126 & 0,144 & 0,044 & \textbf{38,8} \\
bge\_m3 & Superhard & 0,187 & 0,136 & 0,149 & 0,043 & 55,4 \\
jina\_embeddings\_v3 & Superhard & 0,158 & 0,095 & 0,111 & 0,037 & 201,7 \\
\hline
\end{tabular}%
}
\end{table}
\FloatBarrier

\textbf{Вывод.} Лучшее семантическое представление даёт \texttt{e5\_large} (выигрыш по Hit@5 на всех трёх уровнях), поэтому он выбран в оптимальный пайплайн. \texttt{e5\_base} --- компромисс для сценариев с жёстким бюджетом задержки (примерно вдвое быстрее при умеренной потере качества). \texttt{BAAI/bge-m3} уступает обоим на hard и superhard, но остаётся кандидатом для гибридного поиска, поскольку поддерживает разреженные векторы.

\subsection{Эксперименты с разбиением на чанки}

Для chunking-экспериментов использовался baseline dense retrieval при $k=5$ и сопоставление на уровне документа/пересечения. В этом эксперименте менялась только стратегия разбиения документа на фрагменты: token-based, sentence-based, recursive, semantic и parent-child. Остальные параметры retrieval фиксировались, чтобы сравнить именно влияние границ чанков и размера контекста.

Разбиение особенно критично для налоговых текстов: условия применения нормы, исключения, сроки и категории налогоплательщиков могут находиться в соседних абзацах. Слишком короткий чанк может не содержать полного условия, а слишком длинный ухудшает точность поиска и увеличивает шум. Parent-child-подход проверял компромисс, при котором поиск идёт по более точному child-фрагменту, а в контекст можно передать более полный parent-фрагмент. Лучшие варианты различаются по датасетам: на easy лучше всего сработал \texttt{token\_1024\_256}, на hard --- \texttt{parent\_3072\_child\_768}, на superhard --- \texttt{semantic\_1024\_t08}.

\clearpage
\begin{table}[H]
\centering
\caption{Сравнение стратегий chunking при $k=5$}
\label{tab:chunking_results}
\resizebox{\textwidth}{!}{%
\begin{tabular}{llcccc}
\hline
\textbf{Вариант} & \textbf{Датасет} & \textbf{Hit@5} & \textbf{MRR@5} & \textbf{nDCG@5} & \textbf{Precision@5} \\
\hline
parent\_2048\_child\_512 & Easy & 0,503 & 0,385 & 0,407 & 0,162 \\
parent\_2048\_child\_512 & Hard & 0,363 & 0,260 & 0,281 & 0,135 \\
parent\_2048\_child\_512 & Superhard & 0,392 & 0,298 & 0,315 & 0,147 \\
parent\_3072\_child\_768 & Easy & 0,497 & 0,381 & 0,409 & 0,154 \\
parent\_3072\_child\_768 & Hard & \textbf{0,368} & 0,256 & \textbf{0,281} & 0,128 \\
parent\_3072\_child\_768 & Superhard & 0,398 & 0,295 & 0,319 & 0,129 \\
parent\_4096\_child\_1024 & Easy & 0,503 & 0,375 & 0,404 & 0,148 \\
parent\_4096\_child\_1024 & Hard & 0,345 & 0,260 & 0,279 & 0,129 \\
parent\_4096\_child\_1024 & Superhard & 0,404 & 0,299 & 0,318 & 0,125 \\
recursive\_1024 & Easy & 0,503 & 0,387 & 0,397 & 0,188 \\
recursive\_1024 & Hard & 0,351 & 0,247 & 0,258 & 0,139 \\
recursive\_1024 & Superhard & 0,404 & 0,284 & 0,299 & 0,144 \\
recursive\_1536 & Easy & 0,491 & 0,393 & 0,403 & 0,151 \\
recursive\_1536 & Hard & 0,345 & 0,266 & 0,274 & 0,144 \\
recursive\_1536 & Superhard & 0,409 & 0,307 & 0,317 & 0,139 \\
recursive\_legal\_1536 & Easy & 0,491 & \textbf{0,405} & 0,407 & 0,160 \\
recursive\_legal\_1536 & Hard & 0,351 & 0,269 & 0,275 & 0,138 \\
recursive\_legal\_1536 & Superhard & 0,392 & 0,290 & 0,302 & 0,132 \\
semantic\_1024\_t08 & Easy & 0,497 & 0,388 & 0,393 & 0,185 \\
semantic\_1024\_t08 & Hard & 0,333 & 0,243 & 0,249 & 0,132 \\
semantic\_1024\_t08 & Superhard & \textbf{0,433} & 0,293 & 0,312 & 0,153 \\
semantic\_1536\_t075\_skip1 & Easy & 0,491 & 0,369 & 0,389 & 0,156 \\
semantic\_1536\_t075\_skip1 & Hard & 0,339 & 0,257 & 0,271 & 0,124 \\
semantic\_1536\_t075\_skip1 & Superhard & 0,392 & 0,280 & 0,297 & 0,135 \\
sentence\_1024\_128 & Easy & 0,497 & 0,375 & 0,388 & 0,172 \\
sentence\_1024\_128 & Hard & 0,345 & 0,251 & 0,267 & \textbf{0,150} \\
sentence\_1024\_128 & Superhard & 0,409 & 0,300 & 0,315 & 0,146 \\
sentence\_1536\_192 & Easy & 0,485 & 0,375 & 0,390 & 0,150 \\
sentence\_1536\_192 & Hard & 0,363 & \textbf{0,272} & 0,280 & 0,135 \\
sentence\_1536\_192 & Superhard & 0,392 & 0,294 & 0,312 & 0,136 \\
token\_1024\_256 & Easy & \textbf{0,515} & 0,393 & 0,403 & \textbf{0,199} \\
token\_1024\_256 & Hard & 0,363 & 0,258 & 0,269 & 0,147 \\
token\_1024\_256 & Superhard & 0,421 & 0,293 & 0,312 & \textbf{0,154} \\
token\_512\_64 & Easy & 0,480 & 0,366 & 0,382 & 0,191 \\
token\_512\_64 & Hard & 0,351 & 0,250 & 0,256 & \textbf{0,150} \\
token\_512\_64 & Superhard & 0,386 & 0,269 & 0,276 & 0,139 \\
\hline
\end{tabular}%
}
\end{table}
\FloatBarrier

\textbf{Вывод.} Единого победителя нет --- оптимальная стратегия зависит от уровня сложности (\texttt{token\_1024\_256} на easy, parent-child на hard, semantic на superhard). По совокупности (лучший средний Hit@5 и Precision@5) в оптимальный пайплайн выбран \texttt{token\_1024\_256}; parent-child и semantic перспективны и требуют отдельной проверки на очищенном корпусе и совместно с реранкером.

\subsection{Эксперименты с реранкерами}

Эксперименты с реранкерами проверяют двухэтапный retrieval. На первом этапе dense retriever возвращает расширенный список кандидатов; на втором этапе reranker получает пары «вопрос--чанк» и переупорядочивает кандидатов по более точной cross-encoder-оценке. Варьировались модель реранкера и размер набора кандидатов (\texttt{fetch\_k}), а итоговый список оценивался по \texttt{final\_k}. В таблице приведён наиболее показательный режим \texttt{fetch\_k=50}, \texttt{final\_k=5}.

Такой эксперимент важен для нормативных документов, где несколько найденных фрагментов могут быть тематически похожи, но различаться условиями применения. Dense retriever хорошо отбирает кандидатов быстро, но не всегда правильно сортирует близкие документы. Reranker должен поднять выше фрагмент, который действительно отвечает на вопрос, однако это увеличивает latency. Результаты показывают компромисс между качеством и задержкой: \texttt{BAAI/bge-reranker-v2-m3} даёт сильное качество при умеренной задержке; \texttt{jinaai/jina-reranker-v2-base-multilingual} работает заметно быстрее, но слабее на hard и superhard; \texttt{Qwen3-Reranker-0.6B} и \texttt{mixedbread-large} дают конкурентное качество, но существенно дороже по latency.

\begin{table}[H]
\centering
\caption{Сравнение реранкеров при $fetch\_k=50$, $final\_k=5$}
\label{tab:reranker_results}
\resizebox{\textwidth}{!}{%
\begin{tabular}{llccccc}
\hline
\textbf{Реранкер} & \textbf{Датасет} & \textbf{Hit@5} & \textbf{MRR@5} & \textbf{nDCG@5} & \textbf{Precision@5} & \textbf{Rerank, мс} \\
\hline
gte-multilingual-base & Easy & 0,994 & 0,943 & 0,953 & 0,207 & 740 \\
gte-multilingual-base & Hard & 0,731 & 0,612 & 0,641 & 0,163 & 770 \\
gte-multilingual-base & Superhard & 0,877 & 0,740 & 0,773 & 0,193 & 760 \\
BAAI/bge-reranker-v2-m3 & Easy & \textbf{0,994} & 0,935 & 0,949 & 0,207 & 674 \\
BAAI/bge-reranker-v2-m3 & Hard & 0,784 & 0,665 & 0,694 & \textbf{0,179} & 755 \\
BAAI/bge-reranker-v2-m3 & Superhard & \textbf{0,936} & \textbf{0,853} & \textbf{0,870} & \textbf{0,213} & 727 \\
Qwen3-Reranker-0.6B & Easy & \textbf{0,994} & \textbf{0,951} & \textbf{0,962} & \textbf{0,209} & 4417 \\
Qwen3-Reranker-0.6B & Hard & 0,778 & 0,661 & 0,688 & \textbf{0,179} & 4311 \\
Qwen3-Reranker-0.6B & Superhard & \textbf{0,936} & 0,834 & 0,857 & 0,211 & 4294 \\
jina-reranker-v2-base & Easy & 0,988 & \textbf{0,957} & 0,961 & 0,205 & \textbf{249} \\
jina-reranker-v2-base & Hard & 0,702 & 0,581 & 0,607 & 0,157 & \textbf{286} \\
jina-reranker-v2-base & Superhard & 0,883 & 0,750 & 0,775 & 0,194 & \textbf{274} \\
mxbai-rerank-base-v2 & Easy & 0,982 & 0,937 & 0,946 & 0,204 & 2932 \\
mxbai-rerank-base-v2 & Hard & 0,643 & 0,523 & 0,548 & 0,147 & 2726 \\
mxbai-rerank-base-v2 & Superhard & 0,789 & 0,673 & 0,700 & 0,177 & 2786 \\
mxbai-rerank-large-v2 & Easy & 0,988 & 0,935 & 0,949 & 0,206 & 3608 \\
mxbai-rerank-large-v2 & Hard & \textbf{0,801} & \textbf{0,684} & \textbf{0,709} & \textbf{0,179} & 3704 \\
mxbai-rerank-large-v2 & Superhard & \textbf{0,936} & 0,846 & 0,862 & 0,206 & 3578 \\
\hline
\end{tabular}%
}
\end{table}
\FloatBarrier

\textbf{Вывод.} В оптимальный пайплайн выбран \texttt{BAAI/bge-reranker-v2-m3} как лучший компромисс между качеством и задержкой: он показывает топовый или близкий к топовому Hit@5 на всех трёх датасетах ($0{,}994$/$0{,}784$/$0{,}936$) при умеренной задержке ($\approx$700~мс). Модель \texttt{mxbai-rerank-large-v2} немного выше на hard ($0{,}801$), но примерно в пять раз медленнее; \texttt{Qwen3-Reranker-0.6B} сопоставима по качеству, но ещё дороже по latency; \texttt{jina-reranker-v2-base} самая быстрая ($\approx$270~мс), однако заметно слабее на hard и superhard. Таким образом, выбранная модель даёт наибольший прирост качества ранжирования на единицу задержки.

\subsection{Эксперименты с разреженным и гибридным поиском}
\label{sec:hybrid}

Отдельная серия проверяла, помогает ли добавление лексического (sparse) канала к векторному поиску. Базовый dense-ретривер хорошо работает с переформулированными «своими словами» запросами, но может уступать при точном совпадении аббревиатур, номеров статей и реквизитов, где важно буквальное соответствие \cite{robertson1994okapi}. Сравнивались три режима на одном и том же корпусе (коллекция \texttt{rag\_chunks}, 6167 чанков) и при идентичных условиях оценки (\texttt{match\_mode=strict}, попадание на уровне чанка): \texttt{dense}~--- векторный поиск \texttt{BAAI/bge-m3} (воспроизводит baseline из основной матрицы), \texttt{bm25}~--- разреженный поиск Okapi BM25 по токенизированному тексту, \texttt{hybrid}~--- объединение dense и BM25 (по 50 кандидатов каждый) через Reciprocal Rank Fusion (RRF, $k=60$). Реранкер во всех трёх режимах не использовался. Результаты при $k=5$ приведены в таблице~\ref{tab:hybrid_results}; режим \texttt{dense} с точностью до последнего знака совпал с baseline из таблиц~\ref{tab:retrieval_easy}--\ref{tab:retrieval_superhard}, что подтверждает сопоставимость измерений.

\begin{table}[H]
\centering
\caption{Сравнение dense / BM25 / hybrid (RRF) при $k=5$, без реранкера}
\label{tab:hybrid_results}
\resizebox{\textwidth}{!}{%
\begin{tabular}{llcccc}
\hline
\textbf{Датасет} & \textbf{Режим} & \textbf{Hit@5} & \textbf{MRR@5} & \textbf{nDCG@5} & \textbf{Precision@5} \\
\hline
Easy & dense (baseline) & 0,924 & 0,850 & 0,867 & 0,195 \\
Easy & bm25 & \textbf{0,959} & 0,860 & 0,882 & 0,199 \\
Easy & hybrid & \textbf{0,959} & \textbf{0,888} & \textbf{0,903} & \textbf{0,201} \\
\hline
Hard & dense (baseline) & \textbf{0,637} & \textbf{0,521} & \textbf{0,548} & \textbf{0,152} \\
Hard & bm25 & 0,357 & 0,253 & 0,277 & 0,076 \\
Hard & hybrid & 0,556 & 0,407 & 0,442 & 0,125 \\
\hline
Superhard & dense (baseline) & 0,807 & 0,647 & 0,682 & \textbf{0,188} \\
Superhard & bm25 & 0,602 & 0,477 & 0,507 & 0,126 \\
Superhard & hybrid & \textbf{0,819} & \textbf{0,668} & \textbf{0,706} & 0,177 \\
\hline
\end{tabular}%
}
\end{table}
\FloatBarrier

Наблюдения существенно различаются по уровням сложности.
\begin{itemize}
  \item \textbf{Easy.} И BM25, и гибрид превосходят dense: Hit@5 растёт с 0,924 до 0,959, а гибрид дополнительно улучшает ранжирование (MRR@5 0,888 против 0,850, nDCG@5 0,903 против 0,867). Эффект сильнее всего на верхней позиции: Hit@1 у гибрида 0,836 против 0,795 у dense. Это согласуется с гипотезой, что формальные запросы с аббревиатурами (НДФЛ, НДС) и номерами статей выигрывают от точного лексического совпадения.
  \item \textbf{Hard.} Разговорные запросы со строчной буквы и почти без аббревиатур лишены лексических «якорей»: BM25 резко проседает (Hit@5 0,357 против 0,637 у dense), и его подмешивание тянет гибрид вниз (0,556), \emph{ниже} чистого dense. На этом уровне разреженный канал скорее вреден, и лучшим выбором остаётся dense.
  \item \textbf{Superhard.} Несмотря на шум, запросы содержат лексические якоря (ООО, УСН, названия режимов), поэтому гибрид слегка превосходит dense по Hit@5 (0,819 против 0,807), MRR@5 (0,668 против 0,647) и nDCG@5 (0,706 против 0,682), тогда как чистый BM25 остаётся слабым (0,602).
\end{itemize}

Таким образом, равновесное RRF-объединение dense и BM25 полезно, когда запрос содержит точные лексические совпадения (easy, superhard), но ухудшает результат на полностью разговорных запросах (hard), где BM25 добавляет шум. Это указывает на целесообразность не безусловного гибрида, а взвешенного или адаптивного по типу запроса объединения каналов; чистый dense при этом остаётся наиболее устойчивым базовым режимом.

\subsection{Оценка качества генерации}

Для оценки генерации подготовлен отдельный пайплайн, который работает не по голому retrieval-логу, а по JSONL с вопросом, найденным контекстом и уже сгенерированным ответом. Это позволяет запускать ручную разметку и LLM-валидацию независимо: человек может размечать CSV через веб-интерфейс, пока LLM-судья считает оценки. Важно, что LLM-судья не перегенерирует ответы и не смешивает оценку генерации с повторным запуском RAG-пайплайна.

При оценке генерации каждый найденный чанк дополнительно размечается относительно эталона: \texttt{gold\_relevant}, \texttt{hard\_negative} или \texttt{unknown}. Для одного вопроса может существовать несколько \texttt{gold\_relevant} чанков. Это важно для оценки полноты: корректный ответ должен покрывать не только первый найденный фрагмент, но и все существенные условия, сроки и исключения, присутствующие в релевантных чанках, переданных модели.

Ответы генерировались на самом сложном датасете \texttt{eval\_superhard\_dataset.jsonl} (171 вопрос) в нескольких retrieval-конфигурациях: baseline dense и двухэтапная схема с реранкером (с двумя вариантами порядка контекста, см. ниже). Использовалась модель \texttt{qwen/qwen3-14b} в режиме \texttt{no think}; в качестве LLM-судьи~--- \texttt{openai/gpt-oss-20b} с \texttt{reasoning\_effort=low}. Результаты валидации по трём критериям (\texttt{precision}, \texttt{completeness}, \texttt{format}) и согласованность LLM-судьи с ручной разметкой приведены ниже.

Оценка выполняется тремя независимыми валидаторами. Разделение на отдельные критерии нужно для того, чтобы не смешивать фактическую корректность ответа, полноту покрытия релевантного контекста и качество оформления.

\textbf{Precision} оценивает фактическую точность ответа. В этой части precision не означает retrieval Precision@K. Это метрика генерации: проверяется, подтверждаются ли утверждения ответа найденным контекстом, нет ли галлюцинаций, неподтверждённых сроков, сумм, номеров статей, реквизитов писем или обязанностей. Если ответ уверенно даёт налоговый совет при отсутствии релевантного контекста, precision должна быть низкой. Если релевантного контекста нет и модель корректно отказывается отвечать, precision может быть высокой.

\textbf{Completeness} оценивает полноту ответа относительно найденных \texttt{gold\_relevant} чанков. Валидатор проверяет, покрыты ли важные условия, исключения, сроки, категории налогоплательщиков, налоговые режимы, периоды и ограничения. Если релевантный контекст найден, но ответ даёт только общий вывод или игнорирует существенные условия, completeness снижается.

\textbf{Format} оценивает читаемость и профессиональность ответа. Проверяется структура, ясность, отсутствие сырого HTML, сломанной Markdown/LaTeX-разметки, мусорных символов, битой кодировки, грубого или фамильярного тона. Для строгого промпта дополнительно учитывается наличие ожидаемой структуры: краткий вывод, обоснование и источники.

Каждый валидатор выставляет оценку по шкале от 1 до 5. Оценка 5 означает, что ответ полностью соответствует критерию; 1 означает серьёзное нарушение критерия. LLM-судья также сохраняет краткое объяснение оценки и дополнительные диагностические признаки, например факт отказа от ответа, использование \texttt{gold\_relevant} контекста, использование \texttt{hard\_negative} фрагментов как основания, пропуск важной информации или дефекты формата.

\paragraph{Результаты LLM-судьи.}

Сводные оценки судьи по трём retrieval-конфигурациям приведены в таблице~\ref{tab:generation_quality}. Помимо средних баллов \texttt{precision}/\texttt{completeness}/\texttt{format} (шкала 1--5), приведены диагностические доли: \texttt{should\_abstain}~--- доля вопросов, на которых модель \emph{должна} отказаться (нет релевантного контекста), \texttt{did\_abstain}~--- доля фактических отказов, \texttt{missed\_gold}~--- доля ответов, в которых судья зафиксировал пропуск существенной информации из \texttt{gold\_relevant} чанков.

\begin{table}[H]
\centering
\caption{Качество генерации по оценке LLM-судьи (superhard, $n=171$, судья \texttt{gpt-oss-20b})}
\label{tab:generation_quality}
\resizebox{\textwidth}{!}{%
\begin{tabular}{lccccccc}
\hline
\textbf{Конфигурация} & \textbf{precision} & \textbf{completeness} & \textbf{format} & \textbf{avg} & \textbf{should\_abst.} & \textbf{did\_abst.} & \textbf{missed\_gold} \\
\hline
baseline (dense) & \textbf{3,67} & 3,86 & \textbf{5,00} & \textbf{4,18} & 0,187 & 0,064 & \textbf{0,404} \\
reranker, ctx 1-2-3 & 3,60 & \textbf{3,30} & 4,96 & 3,95 & 0,083 & 0,077 & 0,409 \\
reranker, ctx 1-3-2 & 3,60 & 3,19 & 4,95 & 3,92 & 0,066 & 0,060 & 0,488 \\
\hline
\end{tabular}%
}
\end{table}
\FloatBarrier

Ключевые наблюдения. \textbf{Format} стабильно близок к 5,0 во всех конфигурациях: в сохранённых ответах практически нет дефектов оформления (сырого HTML, битой разметки, грубого тона), что ожидаемо при детерминированной генерации с \texttt{temperature=0} и строгом промпте. \textbf{Precision} держится на уровне 3,5--3,7 и слабо зависит от retrieval-метода. Прямое сравнение \textbf{completeness} между baseline и reranker осложнено тем, что валидатор completeness дорабатывался (см. ниже про калибровку), поэтому надёжнее опираться на precision/format, устойчивые к ревизии рубрики, и на сравнения \emph{внутри одной ревизии} (порядок контекста).

\paragraph{Выбор промпта генерации.}

Отдельно сравнивались три системных промпта при фиксированном retrieval (reranker, ctx 1-2-3, superhard, $n=171$): \texttt{default} (нейтральная инструкция отвечать по контексту), \texttt{strict\_citations} (жёсткий запрет внешних знаний, обязательная структура «вывод~--- обоснование~--- источники» и явная фраза отказа при недостатке контекста) и \texttt{strict\_examples} (тот же строгий промпт, дополненный примерами поведения). Результаты приведены в таблице~\ref{tab:generation_prompt}.

\begin{table}[H]
\centering
\caption{Сравнение промптов генерации (reranker, ctx 1-2-3, superhard, $n=171$)}
\label{tab:generation_prompt}
\resizebox{\textwidth}{!}{%
\begin{tabular}{lcccccc}
\hline
\textbf{Промпт} & \textbf{precision} & \textbf{completeness} & \textbf{format} & \textbf{avg} & \textbf{did\_abstain} & \textbf{missing\_struct.} \\
\hline
default & 3,54 & 2,66 & 4,92 & 3,71 & 0,292 & 0,117 \\
strict\_citations & 3,47 & \textbf{3,20} & \textbf{4,96} & 3,88 & 0,071 & \textbf{0,012} \\
strict\_examples & \textbf{3,64} & 3,18 & 4,89 & \textbf{3,91} & 0,089 & 0,041 \\
\hline
\end{tabular}%
}
\end{table}
\FloatBarrier

Нейтральный \texttt{default} оказался худшим: он отказывается отвечать в 29\% случаев (против $\approx$8\% у строгих промптов) и в 11,7\% ответов нарушает требуемую структуру, из-за чего проседают completeness (2,66) и итоговый avg (3,71). Оба строгих промпта существенно лучше и близки между собой: \texttt{strict\_examples} даёт чуть выше precision (3,64) и avg (3,91), тогда как \texttt{strict\_citations} обеспечивает лучшее соблюдение формата (format 4,96, нарушение структуры лишь в 1,2\%) и сбалансированную долю отказов. В качестве промпта по умолчанию выбран \texttt{strict\_citations}: для legal-tech важна предсказуемая структура ответа и корректный отказ, а разница в avg с \texttt{strict\_examples} (3,88 против 3,91) лежит в пределах шума.

\paragraph{Влияние порядка контекста.}

Два варианта реранкера различаются только порядком сборки контекста: \texttt{ctx 1-2-3}~--- естественный порядок реранкера (наиболее релевантный фрагмент первым), \texttt{ctx 1-3-2}~--- перестановка, помещающая более слабый фрагмент в середину (проверка эффекта «lost in the middle» \cite{liu2023lost}). Поскольку обе конфигурации оценивались одной и той же ревизией судьи на одинаковом найденном контексте, их completeness сопоставимы напрямую. Перестановка \texttt{1-3-2} не улучшила, а ухудшила результат: completeness снизилась с 3,30 до 3,19, а доля ответов с пропуском \texttt{gold\_relevant} информации выросла с 0,409 до 0,488; одновременно выросла доля опоры на \texttt{hard\_negative} фрагменты (с 0,095 до 0,120). Таким образом, при текущей длине контекста (top-5, лимит $\approx$18000 символов) искусственная перестановка фрагментов по краям не даёт выигрыша, и естественный порядок реранкера предпочтительнее.

\paragraph{Согласованность LLM-судьи с человеком.}

Для проверки надёжности LLM-судьи 50 примеров baseline-конфигурации были размечены вручную по тем же трём критериям. Сравнение оценок человека и судьи приведено в таблице~\ref{tab:judge_agreement}. Поскольку оценки выставляются по порядковой шкале 1--5, согласие измеряется сразу несколькими взаимодополняющими метриками:
\begin{itemize}
  \item \textbf{Человек} и \textbf{Судья}~--- средние баллы по соответствующему критерию; их разность показывает систематический сдвиг (завышает или занижает судья относительно человека).
  \item \textbf{Exact}~--- доля примеров, где оценка судьи в точности совпала с оценкой человека (строгое попадание балл-в-балл).
  \item \textbf{$\pm 1$}~--- доля примеров, где оценки расходятся не более чем на один балл; это более мягкий критерий, отражающий «практическое» согласие, при котором расхождение в пределах одного деления шкалы считается несущественным.
  \item \textbf{MAE} (mean absolute error)~--- средняя абсолютная разница между оценками судьи и человека в баллах; в отличие от $\kappa$ выражается прямо в единицах шкалы и нагляднее показывает типичную величину ошибки.
  \item \textbf{$\kappa$} (Cohen's $\kappa$)~--- согласие за вычетом случайного совпадения: $\kappa=1$ означает полное согласие, $\kappa=0$~--- согласие на уровне случайного угадывания. Невзвешенный $\kappa$ трактует любое несовпадение одинаково, без учёта его величины. Если все оценки вырождены в одно значение (как у \texttt{format}, где обе стороны почти всегда ставят 5), $\kappa$ не определён.
  \item \textbf{lin $w\kappa$} и \textbf{quad $w\kappa$} (взвешенные $\kappa$)~--- варианты $\kappa$, учитывающие \emph{величину} расхождения: ошибка на 1 балл штрафуется слабее, чем ошибка на 3 балла. Линейная версия наращивает штраф пропорционально расстоянию, квадратичная~--- пропорционально его квадрату (сильнее наказывает большие расхождения и мягче~--- мелкие). Именно взвешенные $\kappa$ наиболее адекватны порядковой шкале и обычно выше невзвешенного.
\end{itemize}

\begin{table}[H]
\centering
\caption{Согласованность LLM-судьи с ручной разметкой (baseline, $n=50$)}
\label{tab:judge_agreement}
\resizebox{\textwidth}{!}{%
\begin{tabular}{lcccccccc}
\hline
\textbf{Критерий} & \textbf{Человек} & \textbf{Судья} & \textbf{Exact} & \textbf{$\pm 1$} & \textbf{MAE} & \textbf{$\kappa$} & \textbf{lin $w\kappa$} & \textbf{quad $w\kappa$} \\
\hline
precision & 3,06 & 3,42 & 0,58 & 0,86 & 0,60 & 0,47 & 0,67 & 0,79 \\
completeness & 2,92 & 3,58 & 0,42 & 0,82 & 0,78 & 0,28 & 0,54 & 0,72 \\
format & 5,00 & 5,00 & 1,00 & 1,00 & 0,00 & --- & --- & --- \\
\hline
macro & --- & --- & 0,67 & 0,89 & 0,46 & 0,38 & 0,60 & 0,76 \\
\hline
\end{tabular}%
}
\end{table}
\FloatBarrier

LLM-судья систематически оценивает ответы выше человека: по precision на $+0{,}36$, по completeness на $+0{,}66$ балла. Несмотря на сдвиг среднего, ранговое согласие по precision умеренно-сильное (quadratic $w\kappa = 0{,}79$, совпадение с допуском $\pm 1$~--- 86\%), что говорит о пригодности судьи для относительного сравнения конфигураций. По completeness согласие слабее ($\kappa = 0{,}28$), а по format оценки вырождены (и человек, и судья почти всегда ставят 5), из-за чего Cohen's $\kappa$ не определён. Слабое согласие именно по completeness и мотивировало доработку соответствующего валидатора.

\paragraph{Калибровка валидатора completeness.}

Исходная рубрика completeness завышала оценку (среднее судьи 3,86 против 2,92 у человека, систематический сдвиг $+0{,}94$). Рубрика дорабатывалась в две итерации; результаты приведены в таблице~\ref{tab:judge_calibration} (среднее судьи по полному набору; MAE и $\kappa$~--- по размеченным человеком примерам).

\begin{table}[H]
\centering
\caption{Калибровка валидатора completeness относительно человека (среднее человека $=2{,}92$)}
\label{tab:judge_calibration}
\begin{tabular}{lcccc}
\hline
\textbf{Ревизия рубрики} & \textbf{Среднее судьи} & \textbf{MAE} & \textbf{Cohen $\kappa$} & \textbf{$n$ (разм.)} \\
\hline
v1 (исходная) & 3,86 & 0,78 & 0,28 & 50 \\
v2 (промежуточная) & 3,69 & 1,35 & 0,10 & 23 \\
v3 (финальная) & 2,94 & 0,73 & 0,24 & 48 \\
\hline
\end{tabular}
\end{table}
\FloatBarrier

Финальная версия v3 практически устранила сдвиг среднего (2,94 против 2,92 у человека) и немного снизила MAE, однако ранговое согласие осталось умеренным ($\kappa = 0{,}24$): часть расхождений между судьёй и человеком не сводится к простому смещению порога. Промежуточная версия v2 оказалась неудачной (рост MAE и падение $\kappa$) и в финальном пайплайне не используется. Меньшие значения $n$ для v2/v3 объясняются тем, что часть ответов судьи не удалось распарсить в строгий JSON. Format и precision к ревизии рубрики устойчивы, поэтому перекалибровывались только при необходимости.

Таким образом, оценка генерации показывает, что система выдаёт ответы с корректным форматом и приемлемой фактической точностью (precision $\approx 3{,}6$), а основной резерв качества~--- полнота покрытия релевантного контекста (completeness) и корректный отказ при недостатке данных. LLM-судья пригоден для относительного сравнения конфигураций, но его абсолютные оценки следует трактовать с поправкой на систематическое завышение относительно человека.

\subsection{Эксперимент: агентный RAG (agentic RAG)}
\label{sec:agentic}

В отличие от одношагового RAG (один запрос $\rightarrow$ top-$k$ $\rightarrow$ ответ), агентный RAG позволяет языковой модели самой решать, что искать, и выполнять несколько последовательных поисков (multi-hop). Цель эксперимента~--- проверить, помогает ли такая многошаговая декомпозиция для сложных налоговых вопросов.

\paragraph{Конфигурация.} В роли агента используется \texttt{openai/gpt-oss-20b} (через OpenAI-совместимый сервер LM Studio) с \texttt{reasoning\_effort=low} и \texttt{temperature=0}. Агенту доступен ровно один инструмент~--- семантический поиск по Qdrant с необязательным фильтром по источнику: \texttt{search(query, source\_code?)}. Поиск~--- dense, \textbf{без реранкера} (только эмбеддер, ради экономии памяти). Эксперимент выполнен на исходной коллекции \texttt{rag\_chunks} (эмбеддер \texttt{BAAI/bge-m3}), поскольку именно на ней построена «золотая» разметка датасетов и single-shot baseline из основной матрицы; это делает сравнение agentic vs single-shot корректным.

\paragraph{Протокол агента.} Используется текстовый JSON-ReAct цикл: на каждом ходу модель возвращает ровно один JSON-объект~--- либо \texttt{\{"action":"search","query":...,"source\_code":?\}}, либо \texttt{\{"action":"answer","content":...\}}. Найденные фрагменты дедуплицируются по тройке \texttt{(source\_code, external\_id, chunk\_i)} и получают сквозную нумерацию \texttt{[n]}; финальный ответ регенерируется из собранного множества фрагментов с тем же строгим промптом \texttt{strict\_citations}. Лимиты бюджета (не более 6 поисков, по 4 фрагмента на поиск, сниппет $\le$600 символов, $\le$24 уникальных фрагмента) удерживают накопленный контекст в пределах $\approx$10--14K токенов; контекст модели поднят до 32K ради запаса под reasoning.

\paragraph{Качественный пример.} Ниже~--- реальная трасса, в которой агент декомпозировал запутанный вопрос и нашёл документ, который одношаговый поиск пропустил:

\begin{quote}\small\ttfamily
Вопрос: «У нас ООО в Екатеринбурге, торгуем автозапчастями; бухгалтер\\
говорит, надо подать письмо Минфина от 19.03.2015 о субсидиях...»\\[4pt]
шаг 1  search(nalog\_letters):\\
\hspace*{2em}"письмо Минфина от 19.03.2015 субсидии автозапчасти ООО"\\
\hspace*{2em}$\rightarrow$ 4 новых фрагмента [1..4]\\
шаг 2  search(minfin\_special):\\
\hspace*{2em}"подать письмо Минфина от 19.03.2015 как оформить"\\
\hspace*{2em}$\rightarrow$ 2 новых фрагмента [5,6]\\
шаг 3  answer (ответ со ссылками [n])\\[4pt]
Итог: single-shot top-5 не нашёл gold-документ; объединение поисков агента~--- нашло.
\end{quote}

\paragraph{Количественное сравнение.} На случайном подмножестве из 40 вопросов superhard (seed~42) сравнивались одношаговый dense-поиск (top-5) и агентный RAG. Метрика~--- \emph{union Hit@k}: попал ли gold-документ в выдачу хотя бы одного из поисков агента (объединение), против Hit@k одношагового поиска. Это \textbf{верхняя граница покрытия} (было ли доказательство извлечено вообще), а не ранжирующая метрика, поэтому напрямую с одношаговым Hit@k как ранжированием она несопоставима. Результаты~--- в таблице~\ref{tab:agentic_results}.

\begin{table}[H]
\centering
\caption{Агентный RAG vs одношаговый dense-поиск (40 вопросов superhard; колонки --- критерии \texttt{doc} и \texttt{strict})}
\label{tab:agentic_results}
\resizebox{\textwidth}{!}{%
\begin{tabular}{lcccccc}
\hline
\textbf{Метод} & \textbf{Hit (doc)} & \textbf{Hit (strict)} & \textbf{ср. поисков} & \textbf{ср. фрагментов} & \textbf{latency ср.} & \textbf{latency p95} \\
\hline
single-shot dense & \textbf{0,750} & \textbf{0,750} & 1 & 5 & $\approx$0,1 с & $\approx$0,1 с \\
agentic (union) & 0,675 & 0,675 & 2,3 & 6,6 & 49,7 с & 70,6 с \\
\hline
\end{tabular}%
}
\end{table}
\FloatBarrier

\paragraph{Выводы и ограничения.} Результат отрицательный и поучительный: агентный RAG, несмотря на то что извлекает \emph{больше} фрагментов (6,6 против 5), покрывает gold-документ \emph{реже} одношагового поиска (union Hit 0,675 против 0,750). Разбор по примерам показывает, почему: агент нашёл документ, пропущенный одношаговым поиском, в 3 случаях, но в 6 случаях, наоборот, потерял его. Причина~--- переформулировка: агент перефразирует исходный (зашумлённый) вопрос в «более чистые» запросы и теряет сигнал, который dense-поиск по исходной формулировке улавливал. При этом цена многошаговости высока: $\approx$50 с на вопрос против $\approx$0,1 с. Совпадение strict и doc (0,675) означает, что на уровне чанка дополнительного выигрыша тоже нет.

Из этого следуют конкретные ограничения и направления. Во-первых, наивная агентная переформулировка с моделью 20B при \texttt{reasoning\_effort=low} не является бесплатным улучшением retrieval и на разговорных superhard-вопросах вредит покрытию. Во-вторых, очевидный фикс~--- всегда включать дословный исходный запрос как один из поисков (чтобы не терять литеральные совпадения), а также пробовать более сильный агент / \texttt{reasoning\_effort=medium}. В-третьих, потенциальная польза агентного подхода может лежать не в покрытии retrieval, а в \emph{синтезе} ответа (декомпозиция, сопоставление норм и разъяснений), что данная метрика не измеряет. Наконец, агент (\texttt{gpt-oss-20b}) отличается от продакшен-генератора (\texttt{qwen3-14b}), поэтому абсолютные значения не следует переносить на основной пайплайн.

\subsection{Итоговый (оптимальный) пайплайн}
\label{sec:optimal}

По результатам всех серий экспериментов собрана рекомендуемая конфигурация, реализованная в Telegram-боте. Для каждого компонента выбран лучший по соответствующей серии вариант:
\begin{itemize}
  \item \textbf{Предобработка корпуса}: профиль \texttt{clean\_no\_boilerplate} --- лучшие Hit@5/MRR/nDCG (таблица~\ref{tab:preprocessing_results}).
  \item \textbf{Разбиение на чанки}: \texttt{token}, размер 1024, перекрытие 256 (\texttt{token\_1024\_256}) --- лучший средний Hit@5 и Precision@5 (таблица~\ref{tab:chunking_results}).
  \item \textbf{Эмбеддер}: \texttt{intfloat/multilingual-e5-large} --- лучший Hit@5 на всех трёх датасетах (таблица~\ref{tab:embedder_results}).
  \item \textbf{Первичный retrieval}: dense без гибрида, поскольку равновесный RRF улучшает не все датасеты и вредит на разговорных запросах (таблица~\ref{tab:hybrid_results}).
  \item \textbf{Реранкер}: \texttt{BAAI/bge-reranker-v2-m3}, \texttt{fetch\_k=50} --- лучший баланс качества и задержки (таблица~\ref{tab:reranker_results}).
  \item \textbf{Сборка контекста}: естественный порядок ctx 1-2-3 (выше completeness, меньше пропусков gold-информации).
  \item \textbf{Преобразование запроса}: не используется (query rewrite и HyDE ухудшают retrieval, таблица~\ref{tab:retrieval_overall}).
  \item \textbf{Генерация}: \texttt{qwen/qwen3-14b}, режим \texttt{no think}, \texttt{temperature=0}, промпт \texttt{strict\_citations} (таблицы~\ref{tab:generation_quality}, \ref{tab:generation_prompt}).
\end{itemize}

Следует подчеркнуть методологическое ограничение этой сборки: серии chunking, embedder и preprocessing измерялись на отдельных переиндексированных коллекциях с мягким критерием \texttt{doc\_overlap}, тогда как retrieval-матрица, реранкер и гибрид~--- на боевой коллекции со строгим критерием \texttt{strict} (попадание на уровне чанка). Поэтому абсолютные приросты компонентов не складываются напрямую, и пайплайн объединяет лучших «в своей серии», но ещё не проверен в едином сквозном прогоне. Подтверждающая сквозная оценка этой конфигурации (переиндексация корпуса под \texttt{e5-large + clean\_no\_boilerplate + token\_1024\_256} и измерение retrieval и генерации на трёх датасетах)~--- ближайший следующий шаг.

\subsection{Ограничения baseline и направления дальнейших экспериментов}

Текущий baseline и экспериментальная постановка имеют следующие ограничения:
\begin{itemize}
  \item единым базовым режимом оставлен dense retrieval: разреженный (BM25) и гибридный (dense+BM25 RRF) режимы оценены отдельно (таблица~\ref{tab:hybrid_results}), но дают устойчивый выигрыш не на всех датасетах --- на hard гибрид уступает чистому dense;
  \item не охвачены региональные нормативные акты (законы субъектов РФ, решения муниципальных образований);
  \item основные retrieval-таблицы используют строгий критерий попадания на уровне чанка (\texttt{strict}), но серии chunking/embedder/preprocessing измерены по более мягкому критерию \texttt{doc\_overlap}, поэтому там метрики могут засчитывать попадание в правильный документ даже когда найден не самый точный чанк; для оценки генерации это компенсируется разметкой найденных чанков как \texttt{gold\_relevant}, \texttt{hard\_negative} и \texttt{unknown};
  \item «золотые» датасеты собраны автоматически и не являются идеальными: hard-набор содержит меньше реальных опечаток, чем предполагалось промптом, superhard-набор частично шаблонен по бизнес-контексту (низкий TTR из-за повторяющихся конструкций), а все три набора построены из одних и тех же чанков с фиксированным \texttt{seed=42} и потому покрывают пересекающийся набор документов. Это ограничивает интерпретацию абсолютных значений метрик, хотя относительное сравнение методов остаётся полезным.
\end{itemize}

В дальнейшем целесообразно проверить следующие направления.

\textbf{Сквозная оценка оптимального пайплайна.} Серии chunking, embedder и preprocessing измерялись на отдельных переиндексированных коллекциях, поэтому их приросты не складываются напрямую (раздел~\ref{sec:optimal}). Ближайший шаг~--- переиндексировать корпус под итоговую конфигурацию (\texttt{e5-large + clean\_no\_boilerplate + token\_1024\_256}) и провести единый сквозной прогон retrieval и генерации на всех трёх датасетах, чтобы подтвердить, что выбранные «лучшие в своей серии» компоненты дают выигрыш и в совокупности.

\textbf{Доменная адаптация моделей retrieval.} Используемые эмбеддер и реранкер обучены на общих многоязычных данных. Перспективно дообучить их на доменных парах «вопрос--релевантный фрагмент», добытых из «золотого набора», писем ФНС/Минфина и статей НК РФ (контрастное обучение для эмбеддера, дообучение cross-encoder-реранкера). Это должно повысить чувствительность к специфике налоговой терминологии, номеров статей и реквизитов писем, где общие модели путают близкие, но юридически различные формулировки.

\textbf{Дальнейшая настройка chunking.} Первичные эксперименты с token-, sentence-, semantic- и parent-child-разбиением уже проведены, но результаты показывают, что оптимальный вариант зависит от сложности запроса. Следующий шаг --- проверить устойчивость лучших стратегий на очищенном корпусе и на расширенном наборе вопросов.

\textbf{Contextual retrieval и late chunking.} Перед индексацией каждого чанка планируется добавлять к нему краткий контекст всего документа, генерируемый LLM. Такой подход (contextual retrieval / late chunking) позволяет чанку «знать» о своём месте в документе, что особенно важно для нормативных текстов, где отдельный абзац вне контекста может не содержать ключевых терминов.

\textbf{Multi-query retrieval.} Для каждого запроса планируется генерировать 3--5 перефразировок через LLM, выполнять поиск по каждой независимо, а затем объединять результаты через RRF. Это позволит компенсировать случаи, когда исходная формулировка семантически далека от текста документа.

\textbf{Agentic RAG.} Агентный подход (декомпозиция вопроса на подзадачи и несколько последовательных поисков) был прототипирован и оценён (раздел~\ref{sec:agentic}). В продакшен-бот он сознательно \emph{не} включён: на разговорных superhard-вопросах наивная агентная переформулировка ухудшает покрытие retrieval относительно одношагового поиска при кратно большей задержке ($\approx$50~с против $\approx$0,1~с на запрос). Если возвращаться к агентному подходу, перспективны: сохранять дословный исходный запрос среди поисков (чтобы не терять литеральные совпадения), пробовать более сильный агент и более высокий \texttt{reasoning\_effort}, а также оценивать его пользу на уровне \emph{синтеза} ответа (сопоставление норм НК РФ, писем Минфина/ФНС и календарных сроков), а не только покрытия retrieval.

\textbf{Расширение корпуса.} Планируется добавление региональных нормативных актов (законы субъектов РФ, решения муниципальных образований по местным налогам), что снизит долю отказов при запросах, специфичных для конкретного региона.

\textbf{Улучшение «золотого набора».} Планируется добавление вопросов с разговорными расшифровками аббревиатур («НДФЛ --- это тот, что с зарплаты?»), расширение охвата налоговых режимов и типов налогоплательщиков (самозанятые на НПД, физлица при продаже и наследовании недвижимости), а также частичная ручная валидация автоматически сгенерированных вопросов.

\newpage
\section{Краткое описание Telegram-бота}
\label{sec:bot}

Практическим результатом работы является прототип, реализованный в виде Telegram-бота~--- пользовательского интерфейса к RAG-системе. Бот построен на асинхронном фреймворке \texttt{aiogram} и выполняет итоговый (оптимальный) пайплайн, описанный в разделе~\ref{sec:optimal}: dense-поиск эмбеддером \texttt{multilingual-e5-large} по коллекции Qdrant, переранжирование \texttt{BAAI/bge-reranker-v2-m3} и генерация ответа моделью \texttt{qwen/qwen3-14b} со строгим промптом \texttt{strict\_citations}.

Основные особенности реализации:
\begin{itemize}
  \item \textbf{Принятие правил.} Перед первым использованием пользователь подтверждает правила (команда \texttt{/accept} или кнопка); согласие сохраняется per-user. Правила фиксируют справочный (неюридический) характер ответов и факт логирования.
  \item \textbf{Один вопрос~--- один ответ.} Каждый запрос обрабатывается независимо; диалоговый режим (multiturn) сознательно не поддерживается, предыдущие сообщения в контекст не включаются. Это упрощает поведение системы и снимает риск дрейфа темы при уточняющих вопросах.
  \item \textbf{Контроль релевантности.} На вопросы, не относящиеся к налоговой тематике, бот вежливо отказывается отвечать и приводит примеры допустимых запросов, не запуская генерацию по нерелевантному контексту.
  \item \textbf{Оформление ответа и источников.} Ответ проходит постобработку: цитаты \texttt{[n]} перенумеровываются по порядку, к каждому источнику подставляется название документа и кликабельная ссылка на оригинал на сайте ведомства; неиспользованные источники из вывода удаляются.
  \item \textbf{Логирование.} Все вопросы, ответы и найденные источники записываются в JSONL для последующего анализа качества и пополнения «золотого набора».
  \item \textbf{Развёртывание и производительность.} Бот и его зависимости (PostgreSQL, Qdrant, сервер \texttt{llama.cpp}) запускаются через \texttt{docker-compose}. Модели эмбеддера и реранкера загружаются в видеопамять при старте (прогрев), чтобы исключить задержку и всплеск нагрузки на первом запросе.
\end{itemize}

Доступные команды: \texttt{/start} (приветствие и правила), \texttt{/rules} (показать правила), \texttt{/accept} (принять правила), \texttt{/help} (описание возможностей и пример запроса).

\newpage
\section{Выводы}
\label{sec:conclusions}

В рамках работы разработан и экспериментально обоснован прототип налогового консультанта для малого и среднего бизнеса на основе подхода Retrieval-Augmented Generation. Реализован полный контур: сбор и периодическое обновление корпуса официальных налоговых документов с учётом метаданных и редакций, индексация в векторном хранилище Qdrant, двухэтапный retrieval с переранжированием и генерация ответа со ссылками на источники. Поставленные цель и задачи выполнены: построен воспроизводимый прототип и проведена серия экспериментов, позволяющая обоснованно выбрать ключевые компоненты архитектуры.

Основные результаты экспериментов:
\begin{itemize}
  \item \textbf{Реранкер~--- главный источник прироста качества retrieval.} Двухэтапная схема (dense + cross-encoder \texttt{bge-reranker-v2-m3}) устойчиво превосходит одношаговый dense-поиск; это наиболее результативное архитектурное решение из проверенных.
  \item \textbf{Преобразование запроса не оправдало себя.} Query rewrite и особенно HyDE на используемой модели ухудшают retrieval: reranker не способен восстановить документ, не попавший в набор кандидатов первичного поиска.
  \item \textbf{Выбор эмбеддера значим.} Лучшие результаты на всех трёх датасетах показала модель \texttt{multilingual-e5-large}, опередившая в том числе исходный \texttt{bge-m3}.
  \item \textbf{Предобработка и разбиение на чанки влияют на качество.} Лучшими в своих сериях оказались профиль очистки \texttt{clean\_no\_boilerplate} и токенное разбиение \texttt{token\_1024\_256}.
  \item \textbf{Гибридный (dense+BM25) поиск даёт неоднозначный эффект:} помогает на формальных и зашумлённых запросах, но проигрывает чистому dense на разговорных, поэтому в продуктовый пайплайн не включён.
  \item \textbf{Агентный (multi-hop) RAG дал отрицательный результат:} на разговорных вопросах наивная переформулировка ухудшает покрытие относительно одношагового поиска при кратно большей задержке, поэтому в продукт он не вошёл.
  \item \textbf{Качество генерации.} При строгом промпте \texttt{strict\_citations} и \texttt{temperature=0} ответы стабильно корректны по формату и приемлемы по фактической точности; основной резерв качества~--- полнота покрытия релевантного контекста и корректный отказ при недостатке данных. LLM-судья пригоден для относительного сравнения конфигураций, но систематически завышает оценки относительно человека, что учтено при калибровке.
\end{itemize}

Совокупность этих результатов сведена в итоговый (оптимальный) пайплайн (раздел~\ref{sec:optimal}), реализованный в Telegram-боте (раздел~\ref{sec:bot}). Полученное решение подтверждает применимость RAG-подхода для справочной поддержки по налоговым вопросам: система формирует обоснованные ответы со ссылками на первоисточники и корректно отказывается отвечать при недостатке данных, что критично для legal-tech домена. Дальнейшие направления развития (доменная адаптация моделей retrieval, сквозная оценка пайплайна, contextual retrieval, multi-query, расширение корпуса и улучшение «золотого набора») описаны в соответствующем подразделе.

\newpage
\section{Описание применения генеративных моделей}

В ходе работы над проектом генеративные модели применялись как вспомогательный инструмент. Всё сгенерированное содержимое (код, текст, списки источников) проверялось, редактировалось и принималось вручную. Ниже указаны цели применения, использованные модели и сервисы с адресами в сети «Интернет», а также способ их применения.

\textbf{Цели применения.}
\begin{itemize}
  \item сбор и систематизация официальных источников налоговой информации и формирование списка источников для парсинга;
  \item дополнительный фактчекинг~--- проверка отдельных утверждений и формулировок;
  \item написание программного кода проекта (около 90\%): пайплайн сбора и индексации документов, RAG-пайплайн, серии экспериментов, Telegram-бот, скрипты оценки;
  \item частичное написание и редактирование текста настоящего отчёта.
\end{itemize}

\textbf{Использованные модели и сервисы.} Названия и адреса в сети «Интернет» приведены в таблице~\ref{tab:genai}.

\begin{table}[H]
\centering
\small
\caption{Использованные генеративные модели и сервисы}
\label{tab:genai}
\begin{tabular}{lll}
\hline
\textbf{Сервис / модель} & \textbf{Разработчик} & \textbf{Адрес в сети «Интернет»} \\
\hline
ChatGPT & OpenAI & \url{https://chatgpt.com} \\
Claude & Anthropic & \url{https://claude.ai} \\
OpenAI Codex & OpenAI & \url{https://openai.com/codex} \\
Claude Code & Anthropic & \url{https://www.anthropic.com/claude-code} \\
\hline
\end{tabular}
\end{table}
\FloatBarrier

\textbf{Способ применения.}
\begin{itemize}
  \item \textbf{Сбор источников и фактчекинг}~--- \texttt{ChatGPT} и \texttt{Claude} в диалоговом (чат) режиме: текстовые запросы на поиск и систематизацию источников и проверку отдельных фактов; результаты сверялись с первоисточниками вручную.
  \item \textbf{Написание кода}~--- \texttt{OpenAI Codex} и \texttt{Claude Code} как агентные инструменты разработки (командная строка): генерация и правка кода по заданию, с ручным ревью, отладкой и запуском экспериментов.
  \item \textbf{Написание текста работы (частично)}~--- \texttt{ChatGPT} и \texttt{Claude}: черновики формулировок, редактирование и структурирование разделов. В частности:
  \begin{itemize}
    \item раздел «Обзор литературы» был первоначально подготовлен с помощью \texttt{ChatGPT}, после чего проверен и исправлен \texttt{Claude} на фактологическую точность (корректность ссылок на первоисточники и формулировок);
    \item оформление таблиц экспериментов в LaTeX выполнено с помощью \texttt{Claude} для ускорения подготовки текста; все числовые значения при этом проверялись вручную по первичным данным;
    \item краткие описания итогов отдельных экспериментов в тексте также подготовлены с помощью \texttt{Claude};
    \item все содержательные выводы по работе написаны вручную, без применения генеративных моделей.
  \end{itemize}
\end{itemize}


\newpage
\section{Список литературы}
\printbibliography[heading=none]

\end{document}
